% संपूर्ण मराठी ग्रंथातील प्रकरण 23: पुनरावर्ती फलने
% हा संपादनयोग्य स्रोतखंड आहे. संकलनासाठी संपूर्ण स्रोतसंग्रहातील
% अक्षररूपे, पूर्वभाग, चिन्हव्याख्या आणि आकृती-स्रोत आवश्यक आहेत.
% मूळ: Open Logic Project; मजकूर आणि रूपांतर CC BY 4.0.
% यंत्रानुवाद, दुरुस्ती आणि तपासणी: OpenAI Codex —
% GPT-5.6 Sol आणि GPT-6 Sol, दोन्ही Ultra effort.
% दुरुस्ती-पुनरावलोकन: GPT-6.1 Sol, Ultra effort.
\chapter{पुनरावर्ती फलने}\label{cmp:rec::chap}
\begin{editorial}
  या जेरेमी अविगॅड यांच्या पुनरावर्ती फलनांवरील टिपा आहेत; रिचर्ड झॅक
  यांनी त्यांत सुधारणा करून त्यांचा विस्तार केला आहे. या प्रकरणात
  खरोखरच काही सराव आहेत आणि विन्यासमीमांसेच्या अंकगणितीकरणावरील
  चर्चेचा आधार देण्यासाठी ते स्वतंत्रपणे समाविष्ट करता येते.
\end{editorial}








\section{परिचय}\label{cmp:rec:int:sec}

संगणनक्षमतेची गणितीय उपपत्ती विकसित करण्यासाठी सर्वप्रथम
संगणनक्षमतेचे एक \emph{प्रतिमान} विकसित करावे लागते. आज आपण
संगणनक्षमता ही संगणक करतात त्या प्रकारची गोष्ट आहे असे मानतो, आणि संगणक
चिन्हांच्या साहाय्याने काम करतात. परंतु संगणनक्षमतेच्या उपपत्तींच्या
विकासाच्या आरंभी संगणनाचे आदर्श उदाहरण \emph{संख्यात्मक} संगणन होते.
गणितज्ञांना संख्या-सैद्धांतिक फलनांमध्ये, म्हणजे संगणित करता येणाऱ्या
$f\colon \Nat^n \to \Nat$ या फलनांमध्ये, नेहमीच रस होता. म्हणून
संगणनक्षमतेच्या उपपत्तीच्या आरंभी अशाच फलनांचा अभ्यास झाला, यात आश्चर्य
नाही. बेरीज, गुणाकार, (नैसर्गिक संख्यांची) घातांक क्रिया यांसारख्या
संगणनक्षम संख्यात्मक फलनांच्या सर्वाधिक परिचित उदाहरणांत एक रोचक समान
वैशिष्ट्य आहे: त्यांची \emph{पुनरावर्ती रीतीने} व्याख्या करता येते.
त्यामुळे पुनरावर्ती व्याख्यांच्या आधारे \emph{संगणनक्षम फलनाची}
सर्वसाधारण व्याख्या करण्याचा प्रयत्न अगदी स्वाभाविक आहे.
संख्या-सैद्धांतिक फलनांची पुनरावर्ती रीतीने व्याख्या करण्याच्या अनेक
संभाव्य मार्गांपैकी येथे एका विशेषतः सोप्या व्याख्या-आकृतिबंधाला
मध्यवर्ती स्थान मिळते: तथाकथित \emph{आदिम पुनरावर्तन}.

संगणनक्षम फलनांबरोबरच संगणनक्षम संच आणि संबंध यांतही आपल्याला रस असू
शकतो. दिलेली संख्या संचाचा घटक आहे की नाही याचे उत्तर आपण
संगणित करू शकत असू, तर तो संच संगणनक्षम असतो; आणि दिलेला क्रमित
घटकसमूह $\tuple{n_1, \dots, n_k}$ हा संबंधाचा घटक आहे
की नाही हे आपण संगणित करू शकत असू तेव्हा आणि केवळ तेव्हाच तो संबंध
संगणनक्षम असतो. संचाचे किंवा संबंधाचे \emph{जातिबोधक फल} विचारात
घेतल्यास संगणनक्षम संच व संबंध यांची चर्चा संगणनक्षम फलनांच्या चर्चेत
समाविष्ट करता येते. त्यामुळे आदिम पुनरावर्ती संबंधांचीही व्याख्या करता
येते; उदाहरणार्थ, ``$m$ ला $n$ ने निःशेष भाग जातो'' हा संबंध आदिम
पुनरावर्ती आहे.

तथापि, आदिम पुनरावर्ती फलने---केवळ आदिम पुनरावर्तन वापरून परिभाषित करता
येणारी फलने---हीच एकमेव संगणनक्षम संख्या-सैद्धांतिक फलने नाहीत. आदिम
पुनरावर्तनाची अनेक सामान्यीकरणे विचारात घेण्यात आली आहेत; परंतु फलने
संगणित करण्याची सर्वाधिक शक्तिशाली आणि व्यापक मान्यता असलेली अतिरिक्त
पद्धत म्हणजे अपरिबद्ध शोध. यातून \emph{आंशिक पुनरावर्ती फलनांची}
व्याख्या मिळते, आणि त्यासंबंधित व्याख्येतून \emph{सामान्य पुनरावर्ती
  फलनांची} व्याख्या मिळते. सामान्य पुनरावर्ती फलने संगणनक्षम आणि
सर्वत्र परिभाषित असतात; आणि ही व्याख्या आंशिक पुनरावर्ती फलनांपैकी
सर्वत्र परिभाषित असणारी नेमकी फलने लक्षणित करते. पुनरावर्ती फलने
संगणनाच्या प्रत्येक इतर प्रतिमानाचे (ट्यूरिंग यंत्रे, लॅम्डा कलन इत्यादींचे)
अनुकरण करू शकतात आणि म्हणून ती संगणनाच्या अनेक मान्य प्रतिमानांपैकी एक
प्रतिमान ठरतात.












\section{आदिम पुनरावर्तन}\label{cmp:rec:pre:sec}

नैसर्गिक संख्यांचे एक वैशिष्ट्य असे आहे की, उत्तरवर्ती क्रिया~$+1$
सांत वेळा लागू करून प्रत्येक नैसर्गिक संख्या~$0$ पासून गाठता येते---कोणतीही
नैसर्गिक संख्या एकतर~$0$ असते किंवा~$0$ च्या उत्तरवर्तीचा \dots{}
उत्तरवर्ती असते. या तथ्याचा उपयोग करणारे फलन~$h\colon\Nat \to \Nat$
विनिर्दिष्ट करण्याचा एक मार्ग असा: (a)~आदान~$0$ साठी $h$ चे मूल्य काय
आहे ते विनिर्दिष्ट करणे, आणि (b)~$h(x)$ चे मूल्य दिले असता $h(x+1)$ चे
मूल्य कसे संगणित करायचे तेही विनिर्दिष्ट करणे. कारण (a) आपल्याला
$h(0)$ काय आहे ते थेट सांगते, त्यामुळे~$0$ साठी $h$ परिभाषित होते.
आता (b) ने दिलेली सूचना $x=0$ साठी वापरून आपण~$h(0)$ वरून
$h(1) = h(0+1)$ संगणित करू शकतो. तीच सूचना $x=1$ साठी वापरून
आपण~$h(1)$ वरून $h(2) = h(1+1)$ संगणित करतो, आणि असेच पुढे.
प्रत्येक नैसर्गिक संख्या~$x$ साठी, आपण अखेरीस $h(x)$ वरून $h(x+1)$
परिभाषित करतो त्या पायरीपर्यंत पोहोचतो; आधारमूल्यापासून अशी पावले
घेतल्यामुळे सर्व $x \in \Nat$ साठी $h(x)$ परिभाषित होते.

उदाहरणार्थ, पुढील दोन समीकरणांनी आपण $h\colon \Nat \to \Nat$ विनिर्दिष्ट
करतो असे समजा:
\begin{align*}
h(0) & =  1\\
h(x+1) & =  2 \cdot h(x)
\end{align*}
गुणाकार कसा करायचा हे आपल्याला आधीच माहीत असेल, तर वरील (a) आणि~(b)
साठी आवश्यक माहिती ही समीकरणे देतात. दुसरे समीकरण क्रमाक्रमाने लागू
केल्यावर आपल्याला पुढील गोष्टी मिळतात:
\begin{align*}
  h(1) & = 2\cdot h(0) = 2,\\
  h(2) & = 2\cdot h(1) = 2\cdot 2,\\
  h(3) & = 2 \cdot h(2) = 2\cdot 2 \cdot 2,\\
  & \vdots
\end{align*}
आपण विनिर्दिष्ट केलेले फलन~$h$ हे $h(x) = 2^x$ आहे, हे दिसते.

नैसर्गिक संख्यांच्या वैशिष्ट्यपूर्ण गुणधर्मामुळे हे दोन निकष पूर्ण करणारे
एकमेव फलन~$h$ असते. अशा समीकरणांच्या जोडीला फलन~$h$ ची
\emph{आदिम पुनरावर्तनाने केलेली व्याख्या} म्हणतात. तिला तसे म्हणण्याचे
कारण म्हणजे आपण~$h$ ची व्याख्या ``पुनरावर्ती रीतीने'' करतो, म्हणजे
व्याख्येत, विशेषतः दुसऱ्या समीकरणाच्या उजव्या बाजूला, $h$ स्वतःच येते.
ती ``आदिम'' असते, कारण $h(x+1)$ परिभाषित करताना आपण केवळ~$h(x)$ चे
मूल्य, म्हणजे अगदी आधीचे मूल्य, वापरतो. फलनाची~$\Nat$ वर पुनरावर्ती
रीतीने व्याख्या करण्याचा हा सर्वांत सोपा मार्ग आहे.

बेरीज आणि गुणाकार यांसारखी आणखी मूलभूत फलनेही आपण आदिम पुनरावर्तनाने
परिभाषित करू शकतो. मात्र या प्रकरणांत संबंधित फलने $2$-स्थानी असतात.
आपण आदानस्थानांपैकी एक स्थिर ठेवतो आणि दुसरे पुनरावर्तनासाठी वापरतो.
उदाहरणार्थ, $\Add(x, y)$ परिभाषित करण्यासाठी आपण~$x$ स्थिर ठेवून आधी
$y=0$ साठी मूल्य परिभाषित करू शकतो आणि नंतर~$y$ च्या साहाय्याने $y+1$
साठीचे मूल्य परिभाषित करू शकतो. $x$ स्थिर असल्यामुळे ते परिभाषक
समीकरणांच्या डाव्या आणि उजव्या दोन्ही बाजूंना येईल.
\begin{align*}
\Add(x,0) & =  x\\
\Add(x,y+1) & =  \Add(x,y)+1
\end{align*}
ही समीकरणे सर्व $x$ \emph{आणि}~$y$ साठी $\Add$ चे मूल्य विनिर्दिष्ट
करतात. उदाहरणार्थ, $\Add(2,3)$ शोधण्यासाठी आपण $x = 2$ साठी परिभाषक
समीकरणे लागू करतो: पहिल्या समीकरणाने $\Add(2,0) = 2$ शोधतो, आणि नंतर
दुसरे समीकरण क्रमाक्रमाने वापरून $\Add(2,1) = 2 + 1 = 3$,
$\Add(2, 2) = 3 + 1 = 4$, $\Add(2, 3) = 4 + 1 = 5$ शोधतो.

$\Add$ च्या व्याख्येत आपण दुसऱ्या समीकरणाच्या उजव्या बाजूला $+$ वापरले,
पण केवळ~$1$ ची भर घालण्यासाठी. दुसऱ्या शब्दांत, आपण उत्तरवर्ती फलन
$\Succ(z) = z+1$ वापरले आणि $\Add(x,y+1)$ परिभाषित करण्यासाठी ते मागील
मूल्य~$\Add(x,y)$ ला लागू केले. म्हणून मागील मूल्याला लागू केलेल्या एका
फलनाच्या रूपात पुनरावर्ती व्याख्येचा विचार करता येतो. परंतु त्या फलनाला
केवळ मागील मूल्यावरच नव्हे, तर $x$ आणि~$y$ यांवरही अवलंबून राहण्याची
मुभा दिल्याने काही बिघडत नाही---आणि कधीकधी ते आवश्यक असते. पुढील
समीकरणे पाहा:
\begin{align*}
  \Mult(x,0) & =  0 \\
  \Mult(x,y+1) & =  \Add(\Mult(x,y),x)
\end{align*}
ही, आधीचे मूल्य $\Mult(x,y)$ आणि पहिले आदान~$x$ या दोहोंना फलन
$\Add$ लागू करून मिळणारी, फलन $\Mult$ ची आदिम पुनरावर्ती व्याख्या आहे.
ती सर्व आदाने $x$ आणि~$y$ यांच्यासाठी फलन~$\Mult(x,y)$ परिभाषितही करते.
उदाहरणार्थ, $\Mult(2,3)$ हे $\Mult(2,0)$, $\Mult(2,1)$,
$\Mult(2,2)$ आणि~$\Mult(2,3)$ क्रमाक्रमाने संगणित करून निर्धारित होते:
\begin{align*}
  \Mult(2,0) & = 0\\
  \Mult(2,1) & = \Mult(2,0+1) =
  \Add(\Mult(2,0), 2) = \Add(0, 2) = 2\\
  \Mult(2,2) & = \Mult(2,1+1) =
  \Add(\Mult(2,1), 2) = \Add(2, 2) = 4\\
  \Mult(2,3) & = \Mult(2,2+1) =
  \Add(\Mult(2,2), 2) = \Add(4, 2) = 6
\end{align*}

मग सर्वसाधारण आकृतिबंध असा आहे: फलन~$h(x_0, \dots, x_{k-1}, y)$ ची
आदिम पुनरावर्ती व्याख्या देण्यासाठी आपण दोन समीकरणे देतो. पहिले समीकरण
फलन~$h$ चा उल्लेख न करता $h(x_0, \dots, x_{k-1}, 0)$ चे मूल्य
परिभाषित करते. दुसरे समीकरण $h(x_0, \dots, x_{k-1}, y)$, इतर आदाने
$x_0$, \dots,~$x_{k-1}$ आणि~$y$ यांच्या साहाय्याने $h(x_0,
\dots, x_{k-1}, y+1)$ चे मूल्य परिभाषित करते. त्या दुसऱ्या समीकरणात
केवळ~$h$ चे अगदी आधीचे मूल्य वापरता येते. या दोन समीकरणांच्या उजव्या
बाजूंनी दिलेल्या क्रियाच स्वतः फलने $f$ आणि~$g$ आहेत असे मानल्यास,
आदिम पुनरावर्तनाने नवीन फलन~$h$ परिभाषित करण्याचा सर्वसाधारण आकृतिबंध
असा आहे:
\begin{align*}
  h(x_0, \dots, x_{k-1}, 0) & = f(x_0, \dots, x_{k-1})\\
  h(x_0, \dots, x_{k-1}, y+1) & =
  g(x_0, \dots, x_{k-1}, y, h(x_0, \dots, x_{k-1}, y))
\end{align*}
$\Add$ च्या बाबतीत $k=1$ आणि $f(x_0) = x_0$ (एकरूपता फलन) असते,
तसेच $g(x_0, y, z) = z + 1$ (आपल्या तिसऱ्या आदानाचा उत्तरवर्ती
परत करणारे $3$-स्थानी फलन) असते:
\begin{align*}
  \Add(x_0, 0) & = f(x_0) = x_0\\
  \Add(x_0, y+1) & = g(x_0, y, \Add(x_0, y)) =
  \Succ(\Add(x_0, y))
\end{align*}
$\Mult$ च्या बाबतीत $f(x_0) = 0$ (नेहमी~$0$ परत करणारे स्थिर फलन)
आणि $g(x_0, y, z) = \Add(z,x_0)$ (आपल्या शेवटच्या व पहिल्या आदानांची
बेरीज परत करणारे $3$-स्थानी फलन) असते:
\begin{align*}
  \Mult(x_0,0) & =  f(x_0) = 0 \\
  \Mult(x_0,y+1) & =  g(x_0, y, \Mult(x_0,y)) =
  \Add(\Mult(x_0,y), x_0)
\end{align*}











\section{संयोजन}\label{cmp:rec:com:sec}

$f$ आणि $g$ ही नैसर्गिक संख्यांची दोन एकस्थानी फलने असतील, तर त्यांचे
संयोजन करता येते: $h(x) = f(g(x))$. मग नवीन फलन~$h(x)$ हे $f$ आणि~$g$
या फलनांपासून \emph{संयोजनाने} परिभाषित होते. याचे एकाहून अधिक आदाने
असलेल्या फलनांसाठी सामान्यीकरण करायचे आहे.

ते करण्याचा एक मार्ग असा: $f$ हे $k$-स्थानी फलन आहे आणि $g_0$,
\dots, $g_{k-1}$ ही सर्व $n$-स्थानी असलेली $k$ फलने आहेत असे समजा.
मग पुढीलप्रमाणे नवीन $n$-स्थानी फलन $h$ परिभाषित करता येते:
\[h(x_0, \dots, x_{n-1}) =
f(g_0(x_0, \dots, x_{n-1}), \dots, g_{k-1}(x_0, \dots, x_{n-1}))
\]
$f$ आणि सर्व $g_i$ संगणनक्षम असतील, तर $h$ देखील संगणनक्षम असते:
$h(x_0, \dots, x_{n-1})$ संगणित करण्यासाठी प्रथम प्रत्येक $i = 0$,
\dots,~$k-1$ साठी $y_i = g_i(x_0, \dots, x_{n-1})$ ही मूल्ये संगणित
करा. नंतर ही मूल्ये $f$ ला आदाने म्हणून देऊन
$h(x_0, \dots, x_{n-1}) = f(y_0, \dots, y_{k-1})$ संगणित करा.

आधीपासून उपलब्ध असलेल्या काही फलनांचा उपयोग करून नवीन फलन संगणित करताना
जे घडते त्याचे हे लक्षणचित्रण अनावश्यकपणे बंधनकारक वाटू शकते. एक कारण
म्हणजे, कधीकधी आपण फलनाची सर्व आदाने वापरत नाही; उदाहरणार्थ,~$\Add$
च्या आदिम पुनरावर्ती व्याख्येत वापरण्यासाठी $g(x, y, z) = \Succ(z)$
परिभाषित केले तेव्हा तसेच झाले. पुढील फलने वापरण्याची आपल्याला मुभा आहे
असे समजा:
\[\Proj{n}{i}(x_0, \dots, x_{n-1}) = x_i
\]
फलनांना~$\Proj{k}{i}$ \emph{प्रक्षेपण} फलने म्हणतात:
$\Proj{n}{i}$ हे $n$-स्थानी फलन आहे. मग $g$ ची व्याख्या अशी देता येते:
\[g(x, y, z) = \Succ(\Proj{3}{2}(x, y, z)).
\]
येथे $f$ ची भूमिका $1$-स्थानी फलन $\Succ$ बजावते, म्हणून $k=1$.
तसेच $g_0$ ची भूमिका बजावणारे एक $3$-स्थानी फलन $\Proj{3}{2}$ आहे.
परिणाम म्हणजे तिसऱ्या आदानाचा उत्तरवर्ती परत करणारे $3$-स्थानी फलन.

प्रक्षेपण फलनांमुळे आदानांचा क्रम बदलून किंवा काही आदाने एकरूप करूनही
नवीन फलने परिभाषित करता येतात. उदाहरणार्थ, फलन $h(x) = \Add(x, x)$ ची
व्याख्या अशी देता येते:
\[h(x_0) = \Add(\Proj{1}{0}(x_0),\Proj{1}{0}(x_0)).
\]
येथे $k=2$, $n=1$, $f(y_0,y_1)$ ची भूमिका $\Add$ बजावते, आणि
$g_0(x_0)$ व $g_1(x_0)$ या दोहोंची भूमिका~$\Proj{1}{0}(x_0)$ हे
एकस्थानी प्रक्षेपण फलन (म्हणजे एकरूपता फलन) बजावते.

$f(y_0, y_1)$ हे फलन आपल्याकडे आधीच असेल, तर फलन
$h(x_0, x_1) = f(x_1, x_0)$ ची व्याख्या अशी देता येते:
\[h(x_0, x_1) = f(\Proj{2}{1}(x_0, x_1),\Proj{2}{0}(x_0, x_1)).
\]
येथे $k=2$, $n = 2$, आणि $g_0$ व $g_1$ यांच्या भूमिका अनुक्रमे
$\Proj{2}{1}$ आणि~$\Proj{2}{0}$ बजावतात.

$g_0$, \dots,~$g_{k-1}$ या सर्वांची स्थानसंख्या~$n$ समान असणे आवश्यक
आहे, याबद्दलही शंका वाटू शकते. (लक्षात ठेवा, फलनाची \emph{स्थानसंख्या}
म्हणजे त्याच्या आदानांची संख्या; $n$-स्थानी फलनाची स्थानसंख्या~$n$
असते.) पण प्रक्षेपण फलने समाविष्ट केल्याने अपेक्षित लवचीकता मिळते.
उदाहरणार्थ, $f$ आणि~$g$ ही $3$-स्थानी फलने आहेत आणि पुढीलप्रमाणे
परिभाषित केलेले $h$ हे $2$-स्थानी फलन आहे असे समजा:
\[h(x,y) = f(x,g(x,x,y),y).
\]
प्रक्षेपण फलने वापरून~$h$ ची व्याख्या पुढीलप्रमाणे पुन्हा लिहिता येते:
\[h(x,y) = f(\Proj{2}{0}(x,y), g(\Proj{2}{0}(x,y), \Proj{2}{0}(x,y),
\Proj{2}{1}(x,y)), \Proj{2}{1}(x,y)).
\]
मग $h$ हे $f$ चे $\Proj{2}{0}$, $l$ आणि $\Proj{2}{1}$ यांच्याशी
केलेले संयोजन आहे, येथे
\[l(x,y) = g(\Proj{2}{0}(x,y),\Proj{2}{0}(x,y),\Proj{2}{1}(x,y)),
\]
म्हणजे $l$ हे $g$ चे $\Proj{2}{0}$, $\Proj{2}{0}$ आणि~$\Proj{2}{1}$
यांच्याशी केलेले संयोजन आहे.











\section{आदिम पुनरावर्ती फलने}\label{cmp:rec:prf:sec}

आदिम पुनरावर्तन आणि संयोजन वापरून आधीच्या फलनांपासून नवीन फलने कशी
परिभाषित करता येतात, ते पुन्हा नोंदवू.

\begin{defn}
  \label{cmp:rec:prf:defn:primitive-recursion} $f$ हे $k$-स्थानी फलन ($k\ge 1$)
  आणि $g$ हे $(k+2)$-स्थानी फलन आहे असे समजा. \emph{$f$ आणि $g$ पासून
  आदिम पुनरावर्तनाने परिभाषित केलेले फलन} म्हणजे पुढील समीकरणांनी परिभाषित
  केलेले $(k+1)$-स्थानी फलन~$h$:
  \begin{align*}
    h(x_0,\dots,x_{k-1},0) & =  f(x_0,\dots,x_{k-1}) \\
    h(x_0,\dots,x_{k-1},y+1) & = g(x_0,\dots,x_{k-1}, y, h(x_0,\dots,x_{k-1}, y))
  \end{align*}
\end{defn}

\begin{defn}
  \label{cmp:rec:prf:defn:composition} $f$ हे $k$-स्थानी फलन आणि $g_0$, \dots,
  $g_{k-1}$ ही सर्व $n$-स्थानी असलेली $k$ फलने आहेत असे समजा.
  \emph{$f$ चे $g_0$, \dots,~$g_{k-1}$ यांच्याशी संयोजन करून परिभाषित
    केलेले फलन} म्हणजे पुढील समीकरणाने परिभाषित केलेले $n$-स्थानी
  फलन~$h$:
  \[  h(x_0, \dots, x_{n-1}) =
  f(g_0(x_0, \dots, x_{n-1}), \dots, g_{k-1}(x_0, \dots, x_{n-1})).
  \]
\end{defn}

$\Succ$ आणि, प्रत्येक नैसर्गिक संख्या $n$ व प्रत्येक $i < n$ साठी,
प्रक्षेपण फलने
\[\Proj{n}{i}(x_0,\dots,x_{n-1}) = x_i,
\]
यांच्याबरोबरच फलन~$\Zero(x) = 0$ चा आपण आदिम पुनरावर्ती फलनांत
समावेश करू.

\begin{defn}
  आदिम पुनरावर्ती फलनांचा संच म्हणजे $\Nat^n$ पासून $\Nat$ कडे जाणाऱ्या
  आणि पुढील कलमांनी विगमनाधारित रीतीने परिभाषित केलेल्या फलनांचा संच:
  \begin{enumerate}
  \item $\Zero$ हे आदिम पुनरावर्ती आहे.
  \item $\Succ$ हे आदिम पुनरावर्ती आहे.
  \item प्रत्येक प्रक्षेपण फलन $\Proj{n}{i}$ हे आदिम पुनरावर्ती आहे.
  \item $f$ हे $k$-स्थानी आदिम पुनरावर्ती फलन आणि $g_0$,
    \dots,~$g_{k-1}$ ही $n$-स्थानी आदिम पुनरावर्ती फलने असतील, तर
    $f$ चे $g_0$, \dots,~$g_{k-1}$ यांच्याशी केलेले संयोजन आदिम
    पुनरावर्ती असते.
  \item $f$ हे $k$-स्थानी आदिम पुनरावर्ती फलन आणि $g$ हे $k+2$-स्थानी
    आदिम पुनरावर्ती फलन असेल, तर $f$ आणि $g$ पासून आदिम पुनरावर्तनाने
    परिभाषित केलेले फलन आदिम पुनरावर्ती असते.
\end{enumerate}
\end{defn}

\begin{explain}
अधिक संक्षेपाने, आदिम पुनरावर्ती फलनांचा संच हा $\Zero$, $\Succ$ आणि
प्रक्षेपण फलने~$\Proj{n}{j}$ यांचा समावेश असलेला आणि संयोजन व आदिम
पुनरावर्तन यांखाली बंद असलेला लघुतम संच आहे.

आदिम पुनरावर्ती फलनांच्या संचाचे वर्णन करण्याचा आणखी एक मार्ग म्हणजे
त्याची ``टप्प्यांच्या'' रूपात व्याख्या करणे. आरंभीच्या फलनांचा---$\Zero$,
$\Succ$ आणि प्रक्षेपणांचा---संच $S_0$ ने दर्शवू. ही टप्पा~$0$ वरील
आदिम पुनरावर्ती फलने आहेत. टप्पा $S_i$ परिभाषित झाल्यावर, आधीच~$S_i$
मध्ये असलेल्या फलनांना संयोजन किंवा आदिम पुनरावर्तन यांपैकी एका क्रियेचा
एकदा उपयोग करून मिळणाऱ्या सर्व फलनांचा संच $S_{i+1}$ असू द्या. मग
\[S = \bigcup_{i \in \Nat} S_i
\]
हा सर्व आदिम पुनरावर्ती फलनांचा संच आहे.
\end{explain}

$\Add$ हे आदिम पुनरावर्ती फलन आहे, याची खात्री करू.

\begin{prop}
  बेरीज फलन $\Add(x,y) = x+y$ हे आदिम पुनरावर्ती आहे.
\end{prop}

\begin{proof}
$\Add$ ची $f$ आणि~$g$ या दोन फलनांच्या साहाय्याने केलेली व
\ref{cmp:rec:prf:defn:primitive-recursion} च्या आकृतिबंधाशी जुळणारी आदिम
पुनरावर्ती व्याख्या आपल्याकडे आधीच आहे:
\begin{align*}
  \Add(x_0, 0) & = f(x_0) = x_0\\
  \Add(x_0, y+1) & = g(x_0, y, \Add(x_0, y)) =
  \Succ(\Add(x_0, y))
\end{align*}
म्हणून $f$ आणि $g$ ही आदिम पुनरावर्ती असतील, तर $\Add$ देखील आदिम
पुनरावर्ती असते. $f(x_0) = x_0 = \Proj{1}{0}(x_0)$, आणि प्रक्षेपण
फलने आदिम पुनरावर्ती मानली जातात; म्हणून $f$ हे आदिम पुनरावर्ती आहे.
फलन $g$ हे पुढीलप्रमाणे परिभाषित केलेले तीन-स्थानी फलन $g(x_0, y, z)$ आहे:
\[g(x_0, y, z) = \Succ(z).
\]
यावरून $g$ आदिम पुनरावर्ती आहे असे अजून म्हणता येत नाही, कारण $g$ आणि
$\Succ$ ही अगदी समान फलने नाहीत: $\Succ$ हे एकस्थानी आहे, पण $g$ तीन-स्थानी
असणे आवश्यक आहे. मात्र संयोजनाने आपण $g$ ची ``औपचारिक'' व्याख्या अशी
देऊ शकतो:
\[g(x_0, y, z) = \Succ(\Proj{3}{2}(x_0, y, z))
\]
$\Succ$ आणि $\Proj{3}{2}$ ही आदिम पुनरावर्ती फलने मानली जात असल्यामुळे,
आदिम पुनरावर्ती फलनांपासून संयोजनाने परिभाषित करता येणारे $g$ देखील आदिम
पुनरावर्ती असते.
\end{proof}

\begin{prop}
  \label{cmp:rec:prf:prop:mult-pr}
  गुणाकार फलन $\Mult(x,y) = x \cdot y$ हे आदिम पुनरावर्ती आहे.
\end{prop}

\begin{proof}
  सराव.
\end{proof}

\begin{prob}
  $\Mult$ ची आदिम पुनरावर्ती व्याख्या
  \ref{cmp:rec:prf:defn:primitive-recursion} ने आवश्यक केलेल्या
  रूपात मांडता येते हे दाखवून आणि संबंधित $f$ व~$g$ ही फलने आदिम
  पुनरावर्ती आहेत हे दाखवून \ref{cmp:rec:prf:prop:mult-pr} सिद्ध करा.
\end{prob}

\begin{ex}
आदिम पुनरावर्ती व्याख्येचे आपले अगदी पहिले उदाहरण पुढीलप्रमाणे होते:
\begin{align*}
h(0) & =  1 \\
h(y+1) & =  2 \cdot h(y).
\end{align*}
हे फलन \ref{cmp:rec:prf:defn:primitive-recursion} ने आवश्यक केलेल्या रूपात बसू शकत
नाही, कारण $k=0$. व्याख्येत $1$ आणि $2$ हे स्थिरही येतात. पहिली अडचण
टाळण्यासाठी एक नाममात्र आदान आणू आणि फलन~$h'$ परिभाषित करू:
\begin{align*}
h'(x_0, 0) & =  f(x_0) = 1 \\
h'(x_0, y+1) & =  g(x_0, y, h'(x_0, y)) = 2 \cdot h'(x_0, y).
\end{align*}
फलन $f(x_0) = 1$ हे $\Succ$ आणि $\Zero$ यांपासून संयोजनाने परिभाषित
करता येते: $f(x_0) = \Succ(\Zero(x_0))$. फलन $g$ हे
$g'(z) = 2 \cdot z$ आणि प्रक्षेपणे यांपासून संयोजनाने परिभाषित करता येते:
\begin{align*}
g(x_0, y, z) & = g'(\Proj{3}{2}(x_0, y, z))
\intertext{आणि पुढे $g'$ ची स्वतःची व्याख्या संयोजनाने अशी देता येते:}
g'(z) & = \Mult(g''(z), \Proj{1}{0}(z))
\intertext{आणि}
g''(z) & = \Succ(f(z)),
\end{align*}
येथे $f$ वरीलप्रमाणेच आहे: $f(z) = \Succ(\Zero(z))$. आता आपल्याकडे~$h'$
आल्यावर, पुन्हा संयोजन वापरून $h(y) =
h'(\Proj{1}{0}(y),\Proj{1}{0}(y))$ असे ठेवता येते. यावरून मूलभूत
फलनांपासून संयोजने आणि आदिम पुनरावर्तने यांचा अनुक्रम वापरून $h$ ची
व्याख्या करता येते; म्हणून $h$ हे आदिम पुनरावर्ती आहे.
\end{ex}











\section{आदिम पुनरावर्ती फलनांची चिन्हांकने}\label{cmp:rec:not:sec}

आदिम पुनरावर्ती फलनांचे अचूक विगमनाधारित वर्णन उपलब्ध असल्याचा एक फायदा
म्हणजे त्यांचे वर्णन पद्धतशीरपणे करता येते. उदाहरणार्थ, पुढीलप्रमाणे
प्रत्येक अशा फलनाला एक ``चिन्हांकन'' देता येते. शून्य, उत्तरवर्ती आणि
प्रक्षेपणे यांसाठी अनुक्रमे $\Zero$, $\Succ$ आणि $\Proj{n}{i}$ ही चिन्हे
वापरा. आता $h$ हे $k$-स्थानी फलन~$f$ आणि $n$-स्थानी फलने $g_0$,
\dots,~$g_{k-1}$ यांच्यापासून संयोजनाने परिभाषित केलेले आहे, आणि या
फलनांना अनुक्रमे $F$, $G_0$, \dots,~$G_{k-1}$ ही चिन्हांकने दिली आहेत
असे समजा. मग नवीन चिन्ह $\fn{Comp}_{k,n}$ वापरून फलन $h$ हे
$\fn{Comp}_{k,n}[F,G_0,\dots,G_{k-1}]$ ने दर्शवता येते.

आदिम पुनरावर्तनाने परिभाषित केलेल्या फलनांसाठी समरूप चिन्हांकने वापरता
येतात. $(k+1)$-स्थानी फलन~$h$ हे $k$-स्थानी फलन~$f$ आणि $(k+2)$-स्थानी
फलन~$g$ यांच्यापासून आदिम पुनरावर्तनाने परिभाषित केलेले आहे, आणि $f$
व~$g$ यांना अनुक्रमे $F$ आणि~$G$ ही चिन्हांकने दिली आहेत असे समजा. मग
$h$ ला दिलेले चिन्हांकन $\fn{Rec}_k[F,G]$ असते.

बेरीज फलनाची आदिम पुनरावर्तनाने व्याख्या पुढीलप्रमाणे आहे, हे आठवा:
\begin{align*}
  \Add(x_0, 0) & = \Proj{1}{0}(x_0) = x_0\\
  \Add(x_0, y+1) & = \Succ(\Proj{3}{2}(x_0, y, \Add(x_0, y))) = \Add(x_0, y) +1
\end{align*}
येथे~$f$ ची भूमिका $\Proj{1}{0}$ बजावते आणि~$g$ ची भूमिका
$\Succ(\Proj{3}{2}(x_0, y, z))$ बजावते. हे $1$-स्थानी फलन~$\Succ$
आणि $3$-स्थानी फलन~$\Proj{3}{2}$ यांच्यापासून संयोजनाने परिभाषित
केलेल्या फलनाचा परिणाम असल्यामुळे त्याला
$\fn{Comp}_{1,3}[\Succ,\Proj{3}{2}]$ हे चिन्हांकन दिले जाते. या
व्यवस्थेत बेरीज फलन पुढील चिन्हांकनाने दर्शवता येते:
\[\fn{Rec}_1[\Proj{1}{0},\fn{Comp}_{1,3}[\Succ,\Proj{3}{2}]].
\]
अशी चिन्हांकने कधीकधी उपयोगी ठरतात; उदाहरणार्थ, आदिम पुनरावर्ती फलनांचे
प्रगणन करताना.

\begin{prob}
  $\Mult$ साठी संपूर्ण आदिम पुनरावर्ती चिन्हांकन द्या.
\end{prob}











\section{आदिम पुनरावर्ती फलने संगणनक्षम असतात}\label{cmp:rec:cmp:sec}

फलन $h$ हे आदिम पुनरावर्तनाने परिभाषित केले आहे असे समजा:
\begin{eqnarray*}
h(\vec x, 0)     & = & f(\vec x) \\
h(\vec x, y + 1) & = & g(\vec x, y, h(\vec x, y))
\end{eqnarray*}
आणि $f$ व $g$ ही फलने संगणनक्षम आहेत असे समजा. ($x_0$, \dots,
$x_{k-1}$ यांचे संक्षिप्त रूप म्हणून आपण $\vec x$ वापरतो.) मग
$h(\vec x, 0)$ संगणित करता येते हे उघड आहे, कारण ते फक्त $f(\vec x)$
आहे आणि ते संगणनक्षम आहे असे आपण गृहीत धरले आहे. त्यानंतर
$h(\vec x, 1)$ देखील संगणित करता येते, कारण $1 = 0 + 1$ आणि म्हणून
$h(\vec x, 1)$ चे मूल्य पुढीलप्रमाणे मिळते:
\begin{align*}
 h(\vec x, 1) & = g(\vec x, 0, h(\vec x, 0)) =  g(\vec x, 0, f(\vec x)).
\intertext{याच प्रकारे पुढे जाऊन आपण पुढील मूल्ये संगणित करू शकतो:}
h(\vec x, 2) & = g(\vec x, 1, h(\vec x, 1)) = g(\vec x, 1, g(\vec x, 0, f(\vec x)))\\
h(\vec x, 3) & = g(\vec x, 2, h(\vec x, 2)) = g(\vec x, 2, g(\vec x, 1, g(\vec x, 0, f(\vec x))))\\
h(\vec x, 4) & = g(\vec x, 3, h(\vec x, 3)) = g(\vec x, 3, g(\vec x, 2, g(\vec x, 1, g(\vec x, 0, f(\vec x)))))\\
& \vdots
\end{align*}
अशा रीतीने, सर्वसाधारणपणे $h(\vec x, y)$ संगणित करण्यासाठी
$h(\vec x, 0)$, $h(\vec x, 1)$, \dots, ही मूल्ये क्रमाक्रमाने
$h(\vec x, y)$ पर्यंत संगणित करा.

म्हणून $f$ आणि $g$ ही फलने संगणनक्षम असतील, तर आदिम पुनरावर्ती
व्याख्येपासून नवीन संगणनक्षम फलन मिळते. त्याचप्रमाणे $f$ आणि $g_i$ ही
फलने संगणनक्षम असतील, तर त्यांच्या संयोजनापासूनही संगणनक्षम फलन मिळते.

मूलभूत फलने $\Zero$, $\Succ$ आणि $\Proj{n}{i}$ ही संगणनक्षम आहेत, आणि
संयोजन व आदिम पुनरावर्तन यांमुळे संगणनक्षम फलनांपासून संगणनक्षम फलने
मिळतात. म्हणून प्रत्येक आदिम पुनरावर्ती फलन संगणनक्षम असते.











\section{आदिम पुनरावर्ती फलनांची उदाहरणे}\label{cmp:rec:exa:sec}


आदिम पुनरावर्ती फलनांची काही उदाहरणे आपल्याकडे आधीच आहेत: बेरीज आणि
गुणाकार ही फलने~$\Add$ आणि $\Mult$. एकरूपता फलन $\fn{id}(x) = x$ हे
आदिम पुनरावर्ती आहे, कारण ते फक्त~$\Proj{1}{0}$ आहे. स्थिर फलने
$\fn{const}_n(x) = n$ ही आदिम पुनरावर्ती आहेत, कारण $\Zero$ आणि $\Succ$
यांपासून क्रमिक संयोजनाने ती परिभाषित करता येतात. आदिम पुनरावर्ती
व्याख्यांत स्थिरांचा वापर करायचा असेल तेव्हा हे उपयोगी पडते. उदाहरणार्थ,
फलन $f(x) = 2 \cdot x$ परिभाषित करायचे असेल, तर $\fn{const}_n(x)$ आणि
गुणाकार यांपासून संयोजन करून ते $f(x) =
\Mult(\fn{const}_2(x), \Proj{1}{0}(x))$ असे मिळवता येते. यापुढे आपण
ही युक्ती वापरू.

\begin{prop}
  घातांक फलन $\fn{exp}(x, y) = x^y$ हे आदिम पुनरावर्ती आहे.
\end{prop}

\begin{proof}
  आपण $\fn{exp}$ ची आदिम पुनरावर्ती व्याख्या अशी देऊ शकतो:
  \begin{align*}
    \fn{exp}(x, 0) & = 1\\
    \fn{exp}(x, y+1) & = \Mult(x, \fn{exp}(x,y)).
    \intertext{काटेकोरपणे पाहता ही आदिम पुनरावर्ती फलनांपासून केलेली
      पुनरावर्ती व्याख्या नाही. मात्र औपचारिक रूपात आपल्याकडे पुढील
      व्याख्या आहे:}
    \fn{exp}(x, 0) & = f(x)\\
    \fn{exp}(x, y+1) & = g(x, y, \fn{exp}(x,y)).
    \intertext{येथे}
    f(x) & = \Succ(\Zero(x)) = 1\\
    g(x, y, z) & = \Mult(\Proj{3}{0}(x, y, z), \Proj{3}{2}(x, y, z)) = x \cdot z
  \end{align*}
  आणि म्हणून $f$ व $g$ ही आदिम पुनरावर्ती फलनांपासून संयोजनाने
  परिभाषित केलेली आहेत.
\end{proof}

\begin{prop}
  पुढीलप्रमाणे परिभाषित केलेले पूर्ववर्ती फलन $\fn{pred}(y)$,
  \[  \fn{pred}(y) = \begin{cases}
    0 & \text{जर $y=0$ असेल}\\
    y-1 & \text{अन्यथा}
  \end{cases}
  \]
  हे आदिम पुनरावर्ती आहे.
\end{prop}

\begin{proof}
 पुढील गोष्ट लक्षात घ्या:
 \begin{align*}
   \fn{pred}(0) & = 0 \text{ आणि}\\
   \fn{pred}(y+1) & = y.
 \end{align*}
 ही जवळजवळ आदिम पुनरावर्ती व्याख्या आहे. काटेकोरपणे पाहता ती आदिम
 पुनरावर्तनाने केलेल्या व्याख्येच्या आकृतिबंधात बसत नाही, कारण त्या
 आकृतिबंधासाठी किमान एक अतिरिक्त आदान~$x$ आवश्यक असते. तसेच
 $\fn{pred}(y+1)$ च्या व्याख्येत $\fn{pred}(y)$ प्रत्यक्षात वापरले जात
 नाही, हेही काहीसे वेगळे आहे. पण आपण प्रथम पुढीलप्रमाणे
 $\fn{pred}'(x, y)$ परिभाषित करू शकतो:
 \begin{align*}
   \fn{pred}'(x, 0) & = \Zero(x) = 0,\\
   \fn{pred}'(x, y+1) & = \Proj{3}{1}(x, y, \fn{pred'}(x, y)) = y.
 \end{align*}
आणि नंतर त्यापासून संयोजनाने $\fn{pred}$ परिभाषित करू शकतो; उदाहरणार्थ,
$\fn{pred}(x) = \fn{pred}'(\Zero(x), \Proj{1}{0}(x))$ असे.
\end{proof}

\begin{prop}
  क्रमगुणित फलन $\fn{fac}(x) = \fact{x} = 1 \cdot 2 \cdot 3
  \cdot \dots \cdot x$ हे आदिम पुनरावर्ती आहे.
\end{prop}

\begin{proof}
  त्याची स्पष्ट आदिम पुनरावर्ती व्याख्या अशी आहे:
  \begin{align*}
    \fn{fac}(0) &= 1\\
    \fn{fac}(y+1) & = \fn{fac}(y) \cdot (y+1).
    \intertext{औपचारिक रूपात आपल्याला प्रथम दोन-स्थानी फलन $h$ परिभाषित करावे लागते:}
    h(x, 0) & = \fn{const}_1(x)\\
    h(x, y+1) & = g(x, y, h(x, y))
    \intertext{येथे $g(x, y, z) = \Mult(\Proj{3}{2}(x, y, z),
      \Succ(\Proj{3}{1}(x, y, z)))$; आणि नंतर पुढीलप्रमाणे ठेवावे:}
    \fn{fac}(y) & = h(\Proj{1}{0}(y), \Proj{1}{0}(y)) = h(y,y).
  \end{align*}
  यापुढे आपण थोडी अधिक औपचारिक मोकळीक घेऊ आणि संयोजन व आदिम
  पुनरावर्तन यांच्या साहाय्याने केलेल्या पूर्ण औपचारिक व्याख्या प्रत्येक
  वेळी देणार नाही.
\end{proof}

\begin{prop}
  पुढीलप्रमाणे परिभाषित केलेली छाटलेली वजाबाकी, $x \tsub y$,
  \[  x \tsub y = \begin{cases}
    0 & \text{जर $x < y$ असेल}\\
    x-y & \text{अन्यथा}
  \end{cases}
  \]
  ही आदिम पुनरावर्ती आहे.
\end{prop}

\begin{proof}
  आपल्याकडे पुढील समीकरणे आहेत:
  \begin{align*}
    x \tsub 0 & = x\\
    x \tsub (y+1) & = \fn{pred}(x \tsub y)
  \end{align*}
\end{proof}

\begin{prop}
 $x$ आणि $y$ यांतील अंतर, $\left|x-y\right|$, हे आदिम पुनरावर्ती आहे.
\end{prop}

\begin{proof}
  आपल्याकडे $\left| x-y \right| = (x \tsub y) + (y \tsub x)$ आहे;
  म्हणून अंतर हे $+$ आणि $\tsub$ यांपासून संयोजनाने परिभाषित करता येते,
  आणि ही दोन्ही आदिम पुनरावर्ती आहेत.
\end{proof}

\begin{prop}
  $x$ आणि $y$ यांतील महत्तम, $\fn{max}(x,y)$, हे आदिम पुनरावर्ती आहे.
\end{prop}

\begin{proof}
  $+$ आणि $\tsub$ यांपासून संयोजनाने $\fn{max}(x,y)$ ची व्याख्या अशी
  देता येते:
  \[  \fn{max}(x,y) \defis x + (y \tsub x).
  \]
  $x$ हे महत्तम असेल, म्हणजे $x \ge y$ असेल, तर $y \tsub x = 0$;
  म्हणून $x + (y \tsub x) = x + 0 = x$. $y$ हे महत्तम असेल, तर
  $y \tsub x = y - x$, आणि म्हणून $x + (y \tsub x) = x + (y - x) = y$.
\end{proof}

\begin{prop}
  \label{cmp:rec:exa:prop:min-pr}
  $x$ आणि $y$ यांतील लघुतम, $\fn{min}(x,y)$, हे आदिम पुनरावर्ती आहे.
\end{prop}

\begin{proof}
  सराव.
\end{proof}

\begin{prob}
  \ref{cmp:rec:exa:prop:min-pr} सिद्ध करा.
\end{prob}

\begin{prob}
पुढील फलन आदिम पुनरावर्ती आहे हे दाखवा:
\[f(x, y) =
2^{(2^{\iddots^{2^{x}}})}\raisebox{1ex}{\bigg\rbrace}
\raisebox{1ex}{\text {$y$ वेळा $2$}}\]
\end{prob}

\begin{prob}
पूर्णांक भागाकार $d(x, y) = \lfloor x/y \rfloor$ (म्हणजे दशांश
बिंदूनंतर येणारे सर्व काही वगळलेला भागाकार) हा आदिम पुनरावर्ती आहे, हे
दाखवा. $y = 0$ असताना $d(x, y) = 0$ असे आपण ठरवतो. आदिम पुनरावर्तन
आणि संयोजन वापरून~$d$ ची स्पष्ट व्याख्या द्या.
\end{prob}

\begin{prop}
आदिम पुनरावर्ती फलनांचा संच पुढील दोन क्रियांखाली बंद आहे:
\begin{enumerate}
\item सांत बेरजा: $f(\vec x, z)$ हे आदिम पुनरावर्ती असेल, तर पुढील
फलनही आदिम पुनरावर्ती असते:
\[g(\vec x, y) \defis \sum_{z = 0}^y f(\vec x, z).
\]
\item सांत गुणाकार: $f(\vec x, z)$ हे आदिम पुनरावर्ती असेल, तर पुढील
फलनही आदिम पुनरावर्ती असते:
\[h(\vec x, y) \defis \prod_{z = 0}^y f(\vec x, z).
\]
\end{enumerate}
\end{prop}

\begin{proof}
उदाहरणार्थ, सांत बेरजा पुढील समीकरणांनी पुनरावर्ती रीतीने परिभाषित होतात:
\begin{align*}
  g(\vec x, 0) & = f(\vec x, 0)\\
  g(\vec x, y+1) & = g(\vec x, y) + f(\vec x, y+1).
\end{align*}
\end{proof}












\section{आदिम पुनरावर्ती संबंध}\label{cmp:rec:prr:sec}


\begin{defn}
संबंध $R(\vec x)$ चे जातिबोधक फल
\[\Char{R}(\vec x) = \left\{
  \begin{array}{ll}
  1 & \mbox{जर $R(\vec x)$ असेल} \\
  0 & \mbox{अन्यथा}
  \end{array}
\right.
\]
आदिम पुनरावर्ती असेल, तर त्या संबंधाला आदिम पुनरावर्ती म्हणतात.
\end{defn}

दुसऱ्या शब्दांत, आदिम पुनरावर्ती संबंध $R(\vec x)$ असा उल्लेख केला, तर
$\Char{R}(\vec x) = 1$ या रूपातील संबंध अभिप्रेत असतो; येथे $\Char{R}$
हे कोणत्याही आदानासाठी 1 किंवा 0 यांपैकी एक मूल्य परत करणारे आदिम
पुनरावर्ती फलन असते. उदाहरणार्थ, $x = 0$ असताना आणि केवळ तेव्हाच सत्य
असलेला संबंध $\fn{IsZero}(x)$ हा $\Char{\fn{IsZero}}$ या फलनाशी संबंधित
आहे; हे फलन आदिम पुनरावर्तन वापरून पुढीलप्रमाणे परिभाषित केले जाते:
\begin{align*}
\Char{\fn{IsZero}}(0) & = 1,\\
\Char{\fn{IsZero}}(x+1) & = 0.
\end{align*}

संबंधांचे इतर आदिम पुनरावर्ती फलनांशी संयोजन करता येते, हे स्पष्ट असावे.
म्हणून पुढील संबंधही आदिम पुनरावर्ती आहेत:
\begin{enumerate}
\item समानता संबंध, $x = y$, जो $\fn{IsZero}(\left|x -
  y\right|)$ ने परिभाषित होतो
\item लघुतर अथवा समान संबंध, $x \leq y$, जो $\fn{IsZero}(x
  \tsub y)$ ने परिभाषित होतो
\end{enumerate}


\begin{prop}
  आदिम पुनरावर्ती संबंधांचा संच बूलीय क्रियांखाली बंद आहे; म्हणजे
  $P(\vec x)$ आणि $Q(\vec x)$ आदिम पुनरावर्ती असतील, तर पुढील संबंधही
  आदिम पुनरावर्ती असतात:
  \begin{enumerate}
  \item $\lnot P(\vec x)$
  \item $P(\vec x) \land Q(\vec x)$
  \item $P(\vec x) \lor Q(\vec x)$
  \item $P(\vec x) \lif Q(\vec x)$
  \end{enumerate}
\end{prop}

\begin{proof}
  $P(\vec x)$ आणि $Q(\vec x)$ आदिम पुनरावर्ती आहेत, म्हणजे त्यांची
  जातिबोधक फले $\Char{P}$ आणि $\Char{Q}$ आदिम पुनरावर्ती आहेत, असे
  समजा. $\lnot P(\vec x)$ इत्यादींची जातिबोधक फलेही आदिम पुनरावर्ती
  आहेत हे दाखवायचे आहे.
  \[  \Char{\lnot P}(\vec x) = \begin{cases}
    0 & \text{जर $\Char{P}(\vec x) = 1$ असेल}\\
    1 & \text{अन्यथा}
  \end{cases}
  \]
  आपण $\Char{\lnot P}(\vec x)$ ची व्याख्या $1 \tsub \Char{P}(\vec x)$
  अशी देऊ शकतो.
  \[  \Char{P \land Q}(\vec x) = \begin{cases}
    1 & \text{जर $\Char{P}(\vec x) = \Char{Q}(\vec x) = 1$ असेल}\\
    0 & \text{अन्यथा}
  \end{cases}
  \]
  आपण $\Char{P \land Q}(\vec x)$ ची व्याख्या $\Char{P}(\vec x) \cdot
  \Char{Q}(\vec x)$ किंवा $\fn{min}(\Char{P}(\vec x), \Char{Q}(\vec
  x))$ अशी देऊ शकतो. त्याचप्रमाणे,
  \begin{align*}
    \Char{P \lor Q}(\vec x) & = \fn{max}(\Char{P}(\vec x), \Char{Q}(\vec x)) \text{ आणि}\\
    \Char{P \lif Q}(\vec x) & = \fn{max}(1 \tsub
  \Char{P}(\vec x), \Char{Q}(\vec x)).
  \end{align*}
\end{proof}

\begin{prop}
  आदिम पुनरावर्ती संबंधांचा संच परिबद्ध संख्यकीकरणाखाली बंद आहे; म्हणजे
  $R(\vec x, z)$ हा आदिम पुनरावर्ती संबंध असेल, तर पुढील संबंधही आदिम
  पुनरावर्ती असतात:
  \begin{align*}
    & \bforall{z < y}{R(\vec x, z)} \text{ आणि}\\
    & \bexists{z < y}{R(\vec x, z)}.
  \end{align*}
  $y$ पेक्षा लहान प्रत्येक~$z$ साठी $R(\vec x, z)$ सत्य असेल तेव्हा
  आणि केवळ तेव्हाच $\bforall{z < y}{R(\vec x, z)}$ हे $\vec x$ आणि
  $y$ यांसाठी सत्य असते; $\bexists{z < y}{R(\vec x, z)}$ साठीही समरूप
  अर्थ लागू होतो.
\end{prop}

\begin{proof}
  संकेताप्रमाणे आपण $\bforall{z < 0}{R(\vec x, z)}$ हे सत्य मानतो
  (कारण~$0$ पेक्षा लहान कोणतेही $z$ नसते) आणि
  $\bexists{z < 0}{R(\vec x, z)}$ हे असत्य मानतो. परिबद्ध सार्वत्रिक
  संख्यापक सांत गुणाकार किंवा पुनरावृत्त लघुतमाप्रमाणे कार्य करतो; म्हणजे
  $P(\vec x, y) \defiff \bforall{z < y}{R(\vec x, z)}$ असेल, तर
  $\Char{P}(\vec x, y)$ ची व्याख्या अशी देता येते:
  \begin{align*}
    \Char{P}(\vec x, 0) & = 1\\
    \Char{P}(\vec x, y+1) & =
    \fn{min}(\Char{P}(\vec x, y), \Char{R}(\vec x, y)).
  \end{align*}
  परिबद्ध अस्तित्ववाची संख्यकीकरणाची व्याख्या त्याचप्रमाणे $\fn{max}$
  वापरून देता येते. पर्यायाने, सममूल्यता
  $\bexists{z < y}{R(\vec x, z)}
  \liff \lnot \bforall{z < y}{\lnot R(\vec x, z)}$ वापरून ते परिबद्ध
  सार्वत्रिक संख्यकीकरणापासून परिभाषित करता येते. उदाहरणार्थ,
  $\bexists{x \leq y}{\dots x\dots}$ या रूपातील परिबद्ध संख्यापक
  $\bexists{x < y+1}{\dots x \dots}$ याच्याशी सममूल्य आहे, हे लक्षात घ्या.
\end{proof}

\begin{prob}
  तीन-स्थानी संबंध $x \equiv y \mod n$ (भाजक~$n$ समशेषता) हा आदिम
  पुनरावर्ती आहे, हे दाखवा.
\end{prob}

आणखी एक उपयोगी आदिम पुनरावर्ती फलन म्हणजे पुढीलप्रमाणे परिभाषित केलेले
सशर्त फलन $\fn{cond}(x,y,z)$:
\begin{align*}
  \fn{cond}(x,y,z) & = \begin{cases}
  y & \text{जर $x = 0$ असेल} \\
  z & \text{अन्यथा}.
\end{cases}
\intertext{याची पुनरावर्ती व्याख्या अशी आहे:}
\fn{cond}(0,y,z) & = y,\\
\fn{cond}(x+1,y,z) & = z.
\end{align*}
आदिम पुनरावर्ती संबंधांपासून आदिम पुनरावर्ती फलनांच्या प्रकरणांनुसार
केलेल्या व्याख्यांचे समर्थन करण्यासाठी हे फलन वापरता येते:

\begin{prop}
$g_0(\vec x)$, \dots,~$g_m(\vec x)$ ही आदिम पुनरावर्ती फलने आणि $R_0(\vec
x)$, \dots, $R_{m-1}(\vec x)$ हे आदिम पुनरावर्ती संबंध असतील, तर
पुढीलप्रमाणे परिभाषित केलेले फलन $f$ देखील आदिम पुनरावर्ती असते:
\[f(\vec x) = \begin{cases}
    g_0(\vec x) & \text{जर $R_0(\vec{x})$ असेल} \\
    g_1(\vec x) & \text{जर $R_1(\vec{x})$ असेल आणि $R_0(\vec{x})$ नसेल} \\
    \vdots & \\
    g_{m-1}(\vec x) & \text{जर $R_{m-1}(\vec{x})$ असेल आणि आधीच्या
      कोणत्याही अटी सत्य नसतील}
    \\
    g_m(\vec x) & \mbox{अन्यथा}
\end{cases}
\]
\end{prop}

\begin{proof}
  $m = 1$ असताना हे फक्त पुढीलप्रमाणे परिभाषित केलेले फलन आहे:
  \[  f(\vec x) = \fn{cond}(\Char{\lnot R_0}(\vec x),g_0(\vec x),g_1(\vec
  x)).
  \]
  $m$ हे $1$ पेक्षा मोठे असताना या रूपातील व्याख्यांचे केवळ संयोजन करता
  येते.
\end{proof}










\section{परिबद्ध लघुतमीकरण}\label{cmp:rec:bmi:sec}

\begin{explain}
एखादा गुणधर्म किंवा संबंध~$P$ पूर्ण करणारी लघुतम संख्या म्हणून फलनाची
व्याख्या करणे अनेकदा उपयोगी ठरते. $P$ निर्णेय असेल, तर $P$ पूर्ण करणारी
लघुतम संख्या सापडेपर्यंत $0$, $1$, $2$, \dots, या सर्व संभाव्य संख्यांची
चाचणी घेऊन आपण हे फलन सहज संगणित करू शकतो. या प्रकारच्या अपरिबद्ध
शोधामुळे आपण आदिम पुनरावर्ती फलनांच्या क्षेत्राबाहेर जातो. मात्र स्वतंत्रपणे
दिलेल्या एखाद्या परिबंधापेक्षा \emph{लहान असलेल्या} लघुतम संख्येतच रस
असेल, तर आपण आदिम पुनरावर्ती क्षेत्रात राहतो. दुसऱ्या शब्दांत, आणि थोडे
अधिक सर्वसाधारणपणे, आपल्याकडे आदिम पुनरावर्ती संबंध~$R(x,z)$ आहे असे
समजा. $x$ आणि~$y$ यांना $R(x, z)$ सत्य करणाऱ्या लघुतम $z < y$ कडे
पाठवणारे फलन विचारात घ्या. $R(x, 0)$, $R(x, 1)$, \dots,
$R(x, y-1)$ यांची चाचणी घेऊन तेही संगणित करता येते. पण ते आदिम
पुनरावर्ती का आहे?
\end{explain}

\begin{prop}
$R(\vec x, z)$ आदिम पुनरावर्ती असेल, तर $R(\vec x, z)$ सत्य करणारे
लघुतम~$z$, $y$ पेक्षा लहान असा साक्षीदार असल्यास, परत करणारे आणि अन्यथा
$y$ परत करणारे फलन $m_R(\vec{x}, y)$ देखील आदिम पुनरावर्ती असते.
फलन~$m_R$ आपण असे लिहू:
\[\bmin{z < y}{R(\vec{x}, z)},
\]
\end{prop}

\begin{proof}
$R(\vec{x}, z)$ सत्य करणारे~$z < 0$ असू शकत नाही, कारण मुळातच कोणतेही
$z < 0$ नसते. म्हणून $m_R(\vec x, 0) = 0$.

परिबंध $y + 1$ या रूपाचा असेल, तर तीन प्रकरणे असतात:
\begin{enumerate}
\item $R(\vec{x}, z)$ सत्य करणारे काही $z < y$ असते; अशा वेळी
$m_R(\vec{x}, y+1) = m_R(\vec{x}, y)$.
\item असे कोणतेही~$z<y$ नसते, पण $R(\vec{x}, y)$ सत्य असते; अशा वेळी
$m_R(\vec{x}, y+1) = y$.
\item $R(\vec{x}, z)$ सत्य करणारे कोणतेही $z < y+1$ नसते; अशा वेळी
$m_R(\vec{x}, y+1) = y+1$.
\end{enumerate}
म्हणून आदिम पुनरावर्तनाने $m_R(\vec x, 0)$ ची व्याख्या पुढीलप्रमाणे
देता येते:
\begin{align*}
m_R(\vec{x}, 0) & = 0\\
m_R(\vec{x}, y+1) & =
\begin{cases}
m_R(\vec{x}, y) & \text{जर $m_R(\vec x, y) \neq y$ असेल}\\
y & \text{जर $m_R(\vec x, y) = y$ आणि $R(\vec{x}, y)$ असेल}\\
y+1 & \text{अन्यथा.}
\end{cases}
\end{align*}
$R(\vec x, z)$ सत्य करणारे काही $z<y$ असते तेव्हा आणि केवळ तेव्हाच
$m_R(\vec x, y) \neq y$ असते, हे लक्षात घ्या.
\end{proof}

\begin{prob}
$R(\vec x, z)$ आदिम पुनरावर्ती आहे असे समजा. $R(\vec{x}, z)$ सत्य
करणारे लघुतम~$z$, $y$ पेक्षा लहान असा साक्षीदार असल्यास, परत करणारे आणि
अन्यथा $0$ परत करणारे फलन $m'_R(\vec{x}, y)$ हे~$\Char{R}$ पासून
आदिम पुनरावर्तनाने परिभाषित करा.
\end{prob}











\section{अविभाज्य संख्या}\label{cmp:rec:pri:sec}

परिबद्ध संख्यकीकरण आणि परिबद्ध लघुतमीकरण यांमुळे नेहमी आढळणारी अनेक
फलने व संबंध आदिम पुनरावर्ती आहेत हे दाखवण्यासाठी आपल्याला बरीच साधने
मिळतात. उदाहरणार्थ, ``$x$ ने $y$ ला निःशेष भाग जातो'' हा संबंध विचारात
घ्या; तो $x \mid y$ असा लिहितात. $y$ ला~$x$ ने बाकी न उरता भागता येत
असेल, म्हणजे $y$ ही~$x$ ची पूर्णांक पटीत संख्या असेल, तेव्हा $x \mid y$
हा संबंध सत्य असतो. (ते घडत नसेल, तर आपण $x \nmid y$ असे लिहितो;
भाजक धन असेल तेव्हा याचा अर्थ $y$ ला $x$ ने भागल्यावर बाकी $> 0$ उरते.)
दुसऱ्या शब्दांत, एखाद्या~$z$ साठी $x \cdot z = y$ असेल तेव्हा आणि केवळ
तेव्हाच $x \mid y$. असा~$z$ अस्तित्वात असेल, तर $\leq y$ असलेला एक
साक्षीदार निवडता येतो. म्हणून एखाद्या $z \le y$ साठी $x \cdot z = y$
असेल तेव्हा आणि केवळ तेव्हाच $x \mid y$. त्यामुळे $=$ आणि गुणाकार
यांच्यापासून परिबद्ध अस्तित्ववाची संख्यकीकरण वापरून $x \mid y$ या संबंधाची
व्याख्या पुढीलप्रमाणे करता येते:
\[x \mid y \defiff \bexists{z \leq y}{(x \cdot z) = y}.
\]
अशा रीतीने $x \mid y$ हा संबंध आदिम पुनरावर्ती आहे हे आपण दाखवले.

नैसर्गिक संख्या~$x$ ही $0$ किंवा $1$ नसल्यास आणि तिला फक्त $1$ ने व
स्वतःनेच भाग जात असल्यास तिला \emph{अविभाज्य} म्हणतात. दुसऱ्या शब्दांत,
अविभाज्य संख्यांचे वैशिष्ट्य असे की, जेव्हा $y \mid x$, तेव्हा एकतर
$y = 1$ किंवा~$y=x$. $x$ अविभाज्य आहे की नाही हे तपासण्यासाठी आपल्याला
फक्त सर्व $y \le x$ साठी $y \mid x$ आहे की नाही हे तपासावे लागते, कारण
$y > x$ असल्यास आपोआप~$y \nmid x$. म्हणून, $x$ अविभाज्य असेल तेव्हा आणि
केवळ तेव्हाच सत्य असणाऱ्या $\fn{Prime}(x)$ या संबंधाची व्याख्या
\[\fn{Prime}(x) \defiff x \geq 2 \land \bforall{y \leq x}{(y \mid x \lif y
  = 1 \lor y = x)}
\]
अशी करता येते आणि म्हणून हा संबंध आदिम पुनरावर्ती आहे.

अविभाज्य संख्या $2$, $3$, $5$, $7$, $11$, इत्यादी आहेत. या अनुक्रमात
स्थानांक~$x$ वरची अविभाज्य संख्या परत करणारे $p(x)$ हे फलन विचारात घ्या;
म्हणजे $p(0) = 2$, $p(1) = 3$, $p(2) = 5$, इत्यादी. (सोयीसाठी आपण
$p(x)$ हे अनेकदा $p_x$ असे लिहू; $p_0=2$, $p_1=3$, इत्यादी.)

$x$ पेक्षा मोठी असलेली पहिली अविभाज्य संख्या परत करणारे
$\fn{nextPrime(x)}$ हे फलन आपल्याकडे असेल, तर आदिम पुनरावर्तन वापरून
$p$ ची व्याख्या सहज करता येते:
\begin{align*}
  p(0) & = 2\\
  p(x+1) & = \fn{nextPrime}(p(x))
\end{align*}
$\fn{nextPrime}(x)$ ही $y > x$ आणि $y$ अविभाज्य अशी लघुतम संख्या~$y$ असल्यामुळे
तिचे संगणन अपरिबद्ध शोधाने सहज करता येते. पण यूक्लिडच्या एका निकालामुळे
तिची व्याख्या परिबद्ध लघुतमीकरणानेही करता येते: $x$ पेक्षा मोठी आणि
$\fact{x}+1$ पेक्षा मोठी नसलेली अविभाज्य संख्या नेहमीच अस्तित्वात असते.
\[  \fn{nextPrime(x)} =
  \bmin{y \leq \fact{x}+1}{(y > x \land \fn{Prime}(y))}.
\]
यावरून $\fn{nextPrime}(x)$ आणि म्हणून $p(x)$ ही फलने केवळ संगणनक्षमच नव्हे,
तर आदिम पुनरावर्तीही आहेत.

(उत्सुकता असेल तर यूक्लिडच्या निकालाचा संक्षिप्त पुरावा असा. दोनपेक्षा
लहान मूल्यांसाठी दोन हीच अपेक्षित अविभाज्य संख्या आहे. इतर मूल्यांसाठी
$p_n$ ही $\le x$ असलेली सर्वांत मोठी अविभाज्य संख्या आहे असे समजा आणि
सर्व $\le x$ अविभाज्य संख्यांचा $p =
p_0\cdot p_1 \cdot \dots \cdot p_n$ हा गुणाकार विचारात घ्या. एकतर $p+1$
अविभाज्य असते, किंवा $x$ आणि~$p+1$ यांच्या दरम्यान एखादी अविभाज्य संख्या
असते. ते का? $p+1$ अविभाज्य नाही असे समजा. मग $q \mid p+1$ अशी एखादी
अविभाज्य संख्या असते, जेथे $q < p+1$. $\le x$ असलेल्या कोणत्याही अविभाज्य
संख्येने $p+1$ ला भाग जात नाही. ($p$ च्या व्याख्येनुसार, प्रत्येक
$p_i \le x$ अविभाज्य संख्येने~$p$ ला भाग जातो, म्हणजे बाकी~$0$ उरते.
म्हणून प्रत्येक $p_i \le x$ अविभाज्य संख्येने $p+1$ ला भागल्यास बाकी~$1$
उरते आणि त्यामुळे $p_i \nmid p+1$.) म्हणून $q$ ही $> x$ आणि $< p+1$
असलेली अविभाज्य संख्या आहे. तसेच $p \le \fact{x}$, म्हणून $> x$ आणि $\le
\fact{x}+1$ असलेली अविभाज्य संख्या अस्तित्वात आहे.)

\begin{prob}
परिबद्ध लघुतमीकरण वापरून पूर्णांक भागाकार $d(x, y)$ ची व्याख्या करा.
\end{prob}











\section{क्रमिका}\label{cmp:rec:seq:sec}

आदिम पुनरावर्ती फलनांचा संच उल्लेखनीयरीत्या सुदृढ आहे. तरीही
\emph{क्रमिका} हाताळण्याची पुरेशी पद्धत विकसित केल्यावर आपल्याला आणखी बरेच
काही करता येईल. नैसर्गिक संख्यांच्या सांत क्रमिका आपण पुढील प्रकारे नैसर्गिक
संख्यांशी ओळखू: $\langle a_0, a_1, a_2, \dots, a_k \rangle$ ही क्रमिका
पुढील संख्येशी संगत असेल:
\[p_0^{a_0+1} \cdot p_1^{a_1+1} \cdot p_2^{a_2+1} \cdot \dots \cdot
p_k^{a_k+1}.
\]
उदाहरणार्थ, $\langle 2, 7, 3\rangle$ आणि $\langle 2, 7, 3, 0, 0 \rangle$
या क्रमिकांचे संख्यात्मक संकेतांक वेगळे राहतील याची खात्री करण्यासाठी आपण
घातांकांमध्ये एक मिळवतो. रिकाम्या क्रमिकेचा संकेतांक म्हणून $0$ आणि~$1$
दोन्ही घेता येतात; निश्चिततेसाठी $\emptyseq$ हे~$0$ दर्शवू दे.

क्रमिकांचे हे सांकेतीकरण चालते याचे कारण तथाकथित अंकगणिताचे मूलभूत प्रमेय
होय: प्रत्येक $n \ge 2$ नैसर्गिक संख्या पुढील रूपात एकमेव रीतीने लिहिता
येते:
\[n = p_0^{a_0} \cdot p_1^{a_1} \cdot \dots \cdot p_k^{a_k}
\]
येथे $a_k \ge 1$. त्यामुळे $\tuple{}(a_0,
\dots, a_k) = \tuple{a_0, \dots, a_k}$ हे चित्रण एकास-एक आहे: वेगवेगळ्या
क्रमिका वेगवेगळ्या संख्यांकडे पाठवल्या जातात आणि प्रत्येक संख्येशी जास्तीत
जास्त एकच क्रमिका संगत असते.

आता क्रमिकेची लांबी ठरवणे, तिचा $i$ वा घटक ठरवणे, क्रमिकेच्या शेवटी एक
घटक जोडणे आणि दोन क्रमिकांची जोडणी करणे या सर्व क्रिया आदिम पुनरावर्ती
आहेत हे आपण दाखवू.

\begin{prop}
  क्रमिका~$s$ ची लांबी परत करणारे $\len{s}$ हे फलन आदिम पुनरावर्ती आहे.
\end{prop}

\begin{proof}
  $R(i, s)$ हा संबंध पुढीलप्रमाणे परिभाषित करू:
  \[  R(i, s) \text{ तेव्हा आणि केवळ तेव्हाच }
  p_i \mid s \land p_{i+1} \nmid s.
  \]
  $R$ हा संबंध स्पष्टपणे आदिम पुनरावर्ती आहे. $s$ हा रिकाम्या नसलेल्या
  क्रमिकेचा संकेतांक असेल, म्हणजे
  \[  s = p_0^{a_0+1} \cdot \dots \cdot p_{k}^{a_{k}+1},
  \]
  तेव्हा $p_i \mid s$ असणारी $p_i$ ही सर्वांत मोठी अविभाज्य संख्या असल्यास,
  म्हणजे $i = k$ असल्यास, $R(i,s)$ सत्य असते. म्हणून $s$ ची लांबी $i+1$
  असते तेव्हा आणि केवळ तेव्हाच $p_i$ ही~$s$ ला भाग जाणारी सर्वांत मोठी
  अविभाज्य संख्या असते. त्यामुळे पुढील व्याख्या करता येते:
  \[  \len{s} =
  \begin{cases}
    0 & \text{जर $s = 0$ किंवा $s = 1$ असेल} \\
    1 + \bmin{i < s}{R(i, s)} & \text{अन्यथा.}
  \end{cases}
  \]
  येथे परिबद्ध लघुतमीकरण वापरता येते, कारण $s$ हा क्रमिकेचा संकेतांक असेल
  तेव्हा $R(i,s)$ पूर्ण करणारा फक्त एकच~$i$ असतो आणि असा $i$ अस्तित्वात
  असल्यास तो~$s$ पेक्षा लहान असतो.
\end{proof}

\begin{prop}
  क्रमिका~$s$ च्या शेवटी $a$ जोडल्यावर मिळणारी क्रमिका परत करणारे
  $\fn{append}(s,a)$ हे फलन आदिम पुनरावर्ती आहे.
\end{prop}

\begin{proof}
  $\fn{append}$ ची व्याख्या पुढीलप्रमाणे करता येते:
  \[  \fn{append}(s,a) =
  \begin{cases}
    2^{a+1} & \text{जर $s = 0$ किंवा $s = 1$ असेल} \\
    s \cdot p_{\len{s}}^{a+1} & \text{अन्यथा.}
  \end{cases}
  \]
\end{proof}

\begin{prop}
  $s$ चा $i$ वा घटक परत करणारे $\fn{element}(s,i)$ हे फलन आदिम
  पुनरावर्ती आहे; येथे पहिल्या घटकाला $0$ वा म्हणतात आणि $i$ हा $s$ च्या
  लांबीपेक्षा मोठा किंवा समान असल्यास फलन $0$ परत करते.
\end{prop}

\begin{proof}
$a$ हा~$s$ चा $i$ वा घटक असतो तेव्हा आणि केवळ तेव्हाच $p_i^{a+1}$ ही
$s$ ला भाग जाणारी~$p_i$ ची सर्वांत मोठी घातशक्ती असते, म्हणजे
$p_i^{a+1} \mid s$ पण $p_i^{a+2} \nmid s$. म्हणून:
  \[  \fn{element}(s,i) =
  \begin{cases}
    0 & \mbox{जर $i \geq \len{s}$ असेल} \\
    \bmin{a < s}{(p_i^{a+2} \nmid s)} & \text{अन्यथा.}
  \end{cases}
  \]
\end{proof}

वर परिभाषित केलेल्या फलनांची पूर्ण नावे वापरण्याऐवजी आपण अधिक संक्षिप्त
चिन्हांकन आणू. $\fn{element}(s,i)$ ऐवजी $(s)_i$ वापरू आणि
$\tuple{s_0, \dots, s_k}$ हे पुढील अभिव्यक्तीचे संक्षिप्त रूप मानू:
\[\fn{append}(\fn{append}(\dots \fn{append}(\emptyseq,s_0)
\dots),s_k).
\]
$s$ ची लांबी~$k$ असेल, तर $s$ चे घटक $(s)_0$, \dots,~$(s)_{k-1}$
असतात, हे लक्षात घ्या.

\begin{prop}
दोन क्रमिकांची जोडणी करणारे $\fn{concat}(s,t)$ हे फलन आदिम पुनरावर्ती
आहे.
\end{prop}

\begin{proof}
  आपल्याला पुढील गुणधर्म असलेले $\fn{concat}$ हे फलन हवे आहे:
  \[    \fn{concat}(\tuple{a_0, \dots, a_k}, \tuple{b_0, \dots, b_l}) =
    \tuple{a_0, \dots, a_k, b_0, \dots, b_l}.
  \]
  $t$ चे पहिले $n$ घटक~$s$ च्या शेवटी जोडणारे
  $\fn{hconcat}(s,t,n)$ हे ``साहाय्यक'' फलन वापरू. त्याची आदिम
  पुनरावर्तनाने पुढीलप्रमाणे व्याख्या करता येते:
  \begin{align*}
    \fn{hconcat}(s,t,0) & = s\\
    \fn{hconcat}(s,t,n+1) & = \fn{append}(\fn{hconcat}(s,t,n),(t)_n)
    \intertext{मग $\fn{concat}$ ची व्याख्या पुढीलप्रमाणे करता येते:}
    \fn{concat}(s,t) & = \fn{hconcat}(s,t,\len{t}).
  \end{align*}
\end{proof}

$\fn{concat}(s,t)$ ऐवजी आपण $s \concat t$ असे लिहू.

क्रमिकेच्या लांबीच्या आणि तिच्या सर्वांत मोठ्या घटकाच्या साहाय्याने तिच्या
संख्यात्मक संकेतांकावर परिबंध घालता येणे आपल्याला उपयोगी पडेल. $s$ ही
लांबी~$k$ असलेली क्रमिका आहे आणि तिचा प्रत्येक घटक एखाद्या~$x$ पेक्षा
लहान किंवा समान आहे असे समजा. मग $s$ चे जास्तीत जास्त $k$ अविभाज्य अवयव
असतात; त्यांपैकी प्रत्येक जास्तीत जास्त~$p_{k-1}$ असतो आणि~$s$ च्या
अविभाज्य अवयवीकरणात प्रत्येकाचा घातांक जास्तीत जास्त $x+1$ असतो. दुसऱ्या
शब्दांत, धन लांबीसाठी आपण
\[\fn{sequenceBound}(x,k) = p_{k-1}^{k \cdot (x+1)},
\]
अशी व्याख्या केली, तर वर वर्णन केलेल्या क्रमिका~$s$ चा संख्यात्मक संकेतांक
जास्तीत जास्त~$\fn{sequenceBound}(x,k)$ असतो. शून्य लांबीच्या क्रमिकेसाठी
क्रमिका-परिबंधाचे मूल्य एक मानू.

क्रमिकांवरील असा परिबंध आपल्याला परिबद्ध शोध वापरून नवीन फलने परिभाषित
करण्याची पद्धत देतो. उदाहरणार्थ, परिबद्ध शोध वापरून $\fn{concat}$ साठी
त्याच क्रमिकेचा संकेतांक देणारी पर्यायी व्याख्या करता येते. आपण शोधत असलेल्या वस्तूचे, म्हणजे जोडलेल्या क्रमिकेच्या
संख्यात्मक संकेतांकाचे, आदिम पुनरावर्ती \emph{विनिर्देश} आणि शोध किती दूरवर
करायचा त्यावरील परिबंध एवढेच लिहावे लागते. पुढील व्याख्या हे काम करते:
\begin{align*}
  \fn{concat}(s,t) = {} & \bmin{v < \fn{sequenceBound}(s+t,\len{s} +
    \len{t})}{} \\
  & \quad(\len{v} = \len{s} + \len{t} \land {}\\
    & \qquad (\bforall{i < \len{s}}{((v)_i = (s)_i)}) \land {} \\
      & \qquad (\bforall{j < \len{t}}{((v)_{\len{s}+j} = (t)_j)}))
\end{align*}
या पर्यायी व्याख्येत रिकाम्या क्रमिकेसाठी 0 हा संकेतांक मिळतो. आधीची
पुनरावर्ती व्याख्या दुसरी क्रमिका रिकामी असल्यास पहिल्या क्रमिकेचा
संकेतांक जसाच्या तसा ठेवते. म्हणून दोन्ही क्रमिका रिकाम्या आणि पहिल्या
क्रमिकेचा संकेतांक 1 असल्यास त्या व्याख्येत 1 मिळतो. 0 आणि 1 हे दोन्ही
रिकाम्या क्रमिकेचे संकेतांक असल्यामुळे क्रमिकेचा अर्थ तोच राहतो; मात्र
या एका प्रकरणात संख्यात्मक संकेतांक समान नसतात.


\begin{prob}
पुढील गुणधर्म असलेले आदिम पुनरावर्ती फलन~$\fn{sconcat}(s)$ अस्तित्वात आहे
हे दाखवा:
\[\fn{sconcat}(\tuple{s_0, \dots, s_k}) = s_0 \concat \dots \concat s_k.
\]
\end{prob}

\begin{prob}
पुढील गुणधर्म असलेले आदिम पुनरावर्ती फलन~$\fn{tail}(s)$ अस्तित्वात आहे
हे दाखवा:
\begin{align*}
  \fn{tail}(\emptyseq) & = 0 \text{ आणि}\\
  \fn{tail}(\tuple{s_0, \dots, s_{k}}) & = \tuple{s_1, \dots, s_{k}}.
\end{align*}
\end{prob}

\begin{prop}
  \label{cmp:rec:seq:prop:subseq}
  $s$ च्या $i$ व्या घटकापासून सुरू होणारी आणि लांबी~$n$ असलेली उपक्रमिका
  परत करणारे $\fn{subseq}(s, i, n)$ हे फलन आदिम पुनरावर्ती आहे.
\end{prop}

\begin{proof}
  सराव.
\end{proof}

\begin{prob}
  \ref{cmp:rec:seq:prop:subseq} सिद्ध करा.
\end{prob}











\section{वृक्ष}\label{cmp:rec:tre:sec}

ज्याप्रमाणे क्रमिका संख्यांनी आणि त्यांचे गुणधर्म व त्यांवरील क्रिया त्यांच्या
संकेतांकांवरील आदिम पुनरावर्ती संबंधांनी व फलनांनी दर्शवता येतात, त्याचप्रमाणे
वृक्ष नैसर्गिक संख्यांनी दर्शवणे कधीकधी उपयोगी ठरते. हे करण्यासाठी आपण
क्रमिका आणि त्यांचे संकेतांक वापरू. वृक्ष म्हणजे एकच गाठ (जिला नामांकन असू
शकते), किंवा अनेक उपवृक्षांशी जोडलेली गाठ (जिला नामांकन असू शकते). त्या
गाठीला वृक्षाचे \emph{मूळ} म्हणतात आणि तिला जोडलेल्या उपवृक्षांना त्याचे
\emph{तात्काळ उपवृक्ष} म्हणतात.

वृक्षाचा पुनरावर्ती संकेतांक $\tuple{k, d_1, \dots, d_k}$ या क्रमिकेने
देऊ; येथे $k$ ही तात्काळ उपवृक्षांची संख्या आणि $d_1$, \dots,~$d_k$ हे
त्यांच्या संकेतांकांचे मूल्य आहेत. गाठींना नामांकने असतील, तर ती तात्काळ
उपवृक्षांनंतर समाविष्ट करता येतात. म्हणून फक्त $l$ नामांकन असलेली एक गाठ
मिळून बनलेल्या वृक्षाचा संकेतांक $\tuple{0,l}$ असेल; आणि $l_1$ नामांकन
असलेले मूळ हे $l_2$, $l_3$ नामांकने असलेल्या दोन एक-गाठीच्या वृक्षांना
जोडलेले असेल, तर त्या वृक्षाचा संकेतांक $\tuple{2,
  \tuple{0, l_2}, \tuple{0, l_3}, l_1}$ असेल.

\begin{prop}
  \label{cmp:rec:tre:prop:subtreeseq}
  संकेतांक~$t$ असलेल्या वृक्षाच्या सर्व उपवृक्षांचे संकेतांक घटक म्हणून
  असलेल्या क्रमिकेचा संकेतांक परत करणारे $\fn{SubtreeSeq}(t)$ हे फलन
  आदिम पुनरावर्ती आहे.
\end{prop}

\begin{proof}
  प्रथम, $\fn{ISubtrees}(t) = \fn{subseq}(t, 1, (t)_0)$ हे आदिम पुनरावर्ती
  फलन आहे आणि वृक्ष~$t$ च्या तात्काळ उपवृक्षांचे संकेतांक परत करते, हे लक्षात
  घ्या. आता मूळापासून जास्तीत जास्त $n$ गाठी दूर असलेल्या सर्व उपवृक्षांची,
  संभाव्य पुनरुक्तींसह, क्रमिका काढणारे साहाय्यक फलन
  $\fn{hSubtreeSeq}(t,n)$ परिभाषित करू. मूळापासून जास्तीत जास्त $0$ गाठी
  दूर असलेल्या~$t$ च्या उपवृक्षांची क्रमिका---दुसऱ्या शब्दांत, $t$ च्या
  मुळापासून सुरू होणारी क्रमिका---म्हणजे फक्त~$t$ असलेली क्रमिका. $t$ च्या
  स्तर~$n+1$ पर्यंतच्या सर्व उपवृक्षांची क्रमिका मिळवण्यासाठी स्तर~$n$
  पर्यंतच्या उपवृक्षांची क्रमिका आणि स्तर~$n$ पर्यंतच्या त्या उपवृक्षांच्या
  तात्काळ उपवृक्षांची क्रमिका यांची जोडणी करतो. ही दुसरी क्रमिका मिळवण्यासाठी,
  $f(x)$ हे क्रमिकांचे संकेतांक परत करणारे आदिम पुनरावर्ती फलन असेल, तर
  $g_f(s, k) = f((s)_0) \concat
  \dots \concat f((s)_{k-1})$ हेही आदिम पुनरावर्ती आहे, हे लक्षात घ्या:
    \begin{align*}
      g(s, 0) & = \emptyseq\\
      g(s, k+1) & = g(s, k) \concat f((s)_k)
    \end{align*}
    उदाहरणार्थ, $s$ ही वृक्षांची क्रमिका असेल, तर
    $h(s) = g_{\fn{ISubtrees}}(s, \len{s})$ ही~$s$ च्या घटकांच्या सर्व
    तात्काळ उपवृक्षांची क्रमिका देते. तिच्या साहाय्याने
    $\fn{hSubtreeSeq}$ ची पुढीलप्रमाणे व्याख्या करता येते:
    \begin{align*}
      \fn{hSubtreeSeq}(t, 0) & = \tuple{t} \\
      \fn{hSubtreeSeq}(t, n+1) & = \fn{hSubtreeSeq}(t, n) \concat
      h(\fn{hSubtreeSeq}(t, n)).
    \end{align*}
    संकेतांक~$t$ असलेल्या वृक्षातील उपवृक्षांचा महत्तम स्तर, म्हणजे मूळ आणि
    एखादी पर्णगाठ यांमधील महत्तम अंतर, संकेतांक~$t$ ने परिबद्ध असतो. म्हणून
    संकेतांक~$t$ असलेल्या वृक्षाच्या सर्व उपवृक्षांच्या संकेतांकांची क्रमिका
    $\fn{hSubtreeSeq}(t, t)$ ने मिळते.
\end{proof}

\begin{prob}
  \ref{cmp:rec:tre:prop:subtreeseq} च्या सिद्धतेतील
  $\fn{hSubtreeSeq}$ च्या व्याख्येत सर्वसाधारणपणे पुनरुक्ती येतात. परिणामी
  क्रमिकेत प्रत्येक उपवृक्षाचा संकेतांक फक्त एकदाच येईल याची खात्री करणारी
  पर्यायी व्याख्या द्या.
\end{prob}











\section{पुनरावर्तनाचे इतर प्रकार}\label{cmp:rec:ore:sec}

जोडीकरण आणि क्रमिका वापरून आदिम पुनरावर्तनाच्या अधिक असामान्य (आणि
उपयोगी) रूपांचे समर्थन करता येते. उदाहरणार्थ, पुढील व्याख्येप्रमाणे दोन
फलने एकसामयिकपणे परिभाषित करणे अनेकदा उपयोगी ठरते:
\begin{align*}
h_0(\vec x, 0) & = f_0(\vec x) \\
h_1(\vec x, 0) & = f_1(\vec x) \\
h_0(\vec x, y+1) & = g_0(\vec x, y, h_0(\vec x, y), h_1(\vec x, y)) \\
h_1(\vec x, y+1) & = g_1(\vec x, y, h_0(\vec x, y), h_1(\vec x, y))
\end{align*}
हे \emph{एकसामयिक पुनरावर्तनाचे} एक उदाहरण आहे. फलने परिभाषित करण्याची
आणखी एक उपयोगी पद्धत म्हणजे $h(\vec x, y+1)$ चे मूल्य आधीच्या
$h(\vec x, 0)$, \dots,~$h(\vec x, y)$ या \emph{सर्व} मूल्यांच्या आधारे
देणे; उदाहरणार्थ:
\begin{align*}
h(\vec x, 0) & = f(\vec x) \\
h(\vec x, y+1) & = g(\vec x, y, \tuple{h(\vec x, 0), \dots, h(\vec x, y)}).
\end{align*}
पुढील योजना ही कल्पना अधिक संक्षिप्तपणे व्यक्त करते:
\[h(\vec x, y) = g(\vec x, y, \tuple{h(\vec x, 0), \dots, h(\vec x, y-1)})
\]
येथे $y$ हे $0$ असेल, तेव्हा $g$ चा शेवटचा आदान केवळ रिकामी क्रमिका
असतो, असे समजायचे. दोन्ही मांडण्यांमागील कल्पना अशी की, ``उत्तरवर्ती
टप्प्याचे'' संगणन करताना $h$ ला आतापर्यंत संगणित केलेल्या मूल्यांची संपूर्ण
क्रमिका वापरता येते. याला \emph{मूल्यक्रम पुनरावर्तन} म्हणतात. एका विशिष्ट
उदाहरणात, पुढील प्रकारच्या व्याख्येचे समर्थन करण्यासाठी ते वापरता येते:
\begin{align*}
h(\vec x, y) & = \begin{cases}
  g(\vec x, y, h(\vec x, k(\vec x, y))) & \text{जर $k(\vec x, y) < y$ असेल} \\
  f(\vec x) & \text{अन्यथा.}
\end{cases}
\end{align*}
दुसऱ्या शब्दांत, $y$ येथील $h$ चे मूल्य~$k$ ने दिलेल्या \emph{कोणत्याही}
आधीच्या मूल्यावरील $h$ च्या मूल्याच्या आधारे संगणित करता येते.


\begin{prob}
  शेष फल $r(x,y)$ ची मूल्यक्रम पुनरावर्तनाने व्याख्या करा. ($x$, $y$ या
  नैसर्गिक संख्या असतील आणि $y > 0$ असेल, तर $r(x,y)$ ही~$y$ पेक्षा लहान
  अशी संख्या असते की, एखाद्या~$z$ साठी $x = z\times y + r(x,y)$. निश्चिततेसाठी,
  $y=0$ असल्यास $r(x,0) = 0$ असे मानू.)
\end{prob}

ही फलने नेहमीच्या आदिम पुनरावर्तनाने कशी मिळवायची याचा विचार करा. आदिम
पुनरावर्तनाची आणखी एक शेवटची आवृत्ती अधिक लवचीक आहे: तिच्यात वाटेत
\emph{प्राचल} (बाजूची मूल्ये) बदलण्याची मुभा असते:
\begin{align*}
h(\vec x, 0) & = f(\vec x) \\
h(\vec x, y+1) & = g(\vec x, y, h(k(\vec x), y))
\end{align*}
याचाही नेहमीच्या आदिम पुनरावर्तनाने अनुकार करता येतो. (हे करणे अवघड आहे.
सूचनेसाठी, संगणन हाताने उलट क्रमाने उलगडून पाहा.)











\section{आदिम पुनरावर्ती नसलेली फलने}\label{cmp:rec:npr:sec}

आदिम पुनरावर्ती फलनांमध्ये सहज समजणाऱ्या अर्थाने संगणनक्षम असलेली सर्व
फलने समाविष्ट होत नाहीत. सर्व एकचलीय आदिम पुनरावर्ती फलनांची $f_0$,
$f_1$, $f_2$,~\dots अशी सूची बनवता येते आणि आदान~$y$ वरील $f_x$ चे मूल्य
प्रभावीपणे संगणित करता येते, हे सहज स्पष्ट असावे. दुसऱ्या शब्दांत,
\[g(x,y) = f_x(y)
\]
याने परिभाषित केलेले $g(x,y)$ हे फलन संगणनक्षम आहे. पण मग पुढील फलनही
संगणनक्षम आहे:
\begin{eqnarray*}
h(x) & = & g(x,x) + 1 \\
& = & f_x(x) +1.
\end{eqnarray*}
प्रत्येक आदिम पुनरावर्ती फलन $f_i$ साठी $h$ आणि $f_i$ यांची $i$ येथील
मूल्ये वेगळी असतात. म्हणून $h$ हे संगणनक्षम आहे, पण आदिम पुनरावर्ती नाही;
आणि $g$ विषयीही हेच म्हणता येते. ही कँटर यांच्या कर्णरेषीय युक्तिवादाची
``प्रभावी'' आवृत्ती आहे.

संगणनक्षम पण आदिम पुनरावर्ती नसलेल्या फलनांची अधिक स्पष्ट उदाहरणे देता
येतात. उदाहरणार्थ, $g^n(x)$ हे चिन्हांकन $g(g(\dots g(x)))$ दर्शवू दे;
त्यात एकूण $n$ वेळा $g$ येते. तसेच फलनांची $g_0,g_1,\dots$ ही क्रमिका
पुढीलप्रमाणे परिभाषित करू:
\begin{eqnarray*}
g_0(x) & = & x+1 \\
g_{n + 1}(x) & = & g_n^x(x)
\end{eqnarray*}
प्रत्येक $g_n$ हे फलन आदिम पुनरावर्ती आहे, हे पडताळून पाहता येईल. प्रत्येक
पुढचे फलन त्याच्या आधीच्या फलनापेक्षा खूप अधिक वेगाने वाढते; $g_1(x)$ हे
$2x$ इतके, $g_2(x)$ हे $2^x \cdot x$ इतके आणि $g_3(x)$ साधारणपणे $x$
इतक्या $2$ च्या घातांचा स्तंभ जितक्या वेगाने वाढतो तितक्या वेगाने वाढते.
ॲकरमन--पेटर फलन हे मूलतः $G(x) = g_x(x)$ हे फलन आहे आणि ते कोणत्याही
आदिम पुनरावर्ती फलनापेक्षा अधिक वेगाने वाढते, हे दाखवता येते.

आता आदिम पुनरावर्ती फलनांच्या प्रगणनाकडे परत येऊ. आपण प्रत्येक आदिम
पुनरावर्ती फलनाला सांकेतिक चिन्हांकन दिले आहे, हे लक्षात घ्या; म्हणून
चिन्हांकनांचे प्रगणन करणे पुरेसे आहे. प्रत्येक चिन्हांकन~$F$ ला नैसर्गिक
संख्या $\#(F)$ पुढीलप्रमाणे पुनरावर्ती रीतीने देता येते:
\begin{eqnarray*}
\#(0) & = & \langle 0 \rangle \\
\#(S) & = & \langle 1 \rangle \\
\#(\Proj{n}{i}) & = & \langle 2, n, i \rangle \\
\#(\fn{Comp}_{k,l}[H,G_0,\dots,G_{k-1}]) & = & \langle
3,k,l,\#(H),\#(G_0),\dots,\#(G_{k-1}) \rangle \\
\#(\fn{Rec}_l[G,H]) & = & \langle 4, l, \#(G), \#(H) \rangle
\end{eqnarray*}
येथे मागील विभागातील संकेतांक वापरून संख्यांची प्रत्येक क्रमिका नैसर्गिक
संख्या म्हणून पाहता येते, या तथ्याचा उपयोग केला आहे. परिणामी प्रत्येक
संकेतांकाला नैसर्गिक संख्या मिळते. काही क्रमिका (आणि म्हणून काही संख्या)
चिन्हांकनांशी संगत नसतात, हे उघड आहे. पण $i$ ने अशा चिन्हांकनाला सांकेतिक
केले असेल, तर $i$ ने सांकेतिक केलेल्या त्या चिन्हांकनाचे एकचलीय आदिम
पुनरावर्ती फलन $f_i$ असू दे; आणि
अन्यथा ते स्थिरांक $0$ फलन असू दे. अंतिम परिणाम असा की, एकचलीय आदिम
पुनरावर्ती फलनांचे प्रगणन करण्याची स्पष्ट पद्धत आपल्याला मिळते.

(प्रत्यक्षात, स्थिरांक शून्य फलनासारखी काही फलने सूचीमध्ये एकाहून अधिक वेळा
येतील. हा केवळ आपल्या सांकेतीकरणाचा कृत्रिम परिणाम नाही; स्थिरांक शून्य
फलनाला एकाहून अधिक चिन्हांकने असतात, याचाही तो परिणाम आहे. पुढे आपण पाहू की
या पुनरुक्ती संगणनक्षम रीतीने टाळता येत नाहीत. उदाहरणार्थ, दिलेले चिन्हांकन
स्थिरांक शून्य फलन दर्शवते की नाही हे ठरवणारे कोणतेही संगणनक्षम फलन नसते.)

आता आपण $g(x,y)$ हे फलन $f_x(y)$ ने दिलेले मानू शकतो; येथे $f_x$ म्हणजे
नुकत्याच वर्णन केलेल्या प्रगणनातील फलन. $g(x,y)$ संगणनक्षम आहे हे आपल्याला
कसे माहीत? सहजपणे हे स्पष्ट आहे: $g(x,y)$ संगणित करण्यासाठी प्रथम $x$
``उलगडून'' ते एकचलीय फलनाचे चिन्हांकन आहे की नाही हे पाहा. तसे असल्यास,
आदान~$y$ वरील त्या फलनाचे मूल्य संगणित करा.













\section{आंशिक पुनरावर्ती फलने}\label{cmp:rec:par:sec}


पुनरावर्ती फलनांच्या व्याख्येमागील प्रेरणा समजण्यासाठी पुढील गोष्ट लक्षात
घ्या: संगणनक्षम पण आदिम पुनरावर्ती नसलेली फलने अस्तित्वात आहेत या आपल्या
सिद्धतेतून प्रत्यक्षात आणखी बरेच काही सिद्ध होते. युक्तिवाद साधा होता:
$f_0,f_1,\dots$ या फलनांचे असे प्रगणन करता येते की, $x$ आणि $y$ यांचे फलन
म्हणून $f_x(y)$ संगणनक्षम असते, एवढेच तथ्य आपण वापरले. म्हणून हा युक्तिवाद
\emph{अशा प्रकारे प्रगणन करता येणाऱ्या फलनांच्या कोणत्याही वर्गाला} लागू
होतो. यामुळे आपण अडचणीत येतो: संगणनक्षम फलनांचे स्पष्ट वर्णन आपल्याला हवे
आहे; पण संगणनक्षम फलनांच्या संग्रहाचे अशा प्रकारचे कोणतेही प्रभावी स्पष्ट
वर्णन सर्वसमावेशक असू शकत नाही!

यातून बाहेर पडण्याचा मार्ग म्हणजे \emph{अंशतः} फलनांना स्थान देणे. अंशतः
संगणनक्षम फलनांचे प्रगणन करणे \emph{शक्य आहे}, हे आपण पाहू.
 पुढे आपण आपल्या
कर्णरेषीय युक्तिवादाकडे परत येऊ आणि अंशतः फलने समाविष्ट केल्यावर तो का
चालत नाही याचा शोध घेऊ.

आता प्रश्न असा आहे: सर्व आंशिक पुनरावर्ती फलने मिळवण्यासाठी आदिम पुनरावर्ती
फलनांमध्ये काय भर घालावी लागेल? आपल्याला दोन गोष्टी कराव्या लागतील:
\begin{enumerate}
\item आदिम पुनरावर्ती फलनांची आपली व्याख्या बदलून त्यात अंशतः फलनांनाही
  स्थान द्यावे.
\item काही नवीन अंशतः फलने समाविष्ट होतील अशी काही गोष्ट व्याख्येत
  \emph{जोडावी}.
\end{enumerate}

पहिली गोष्ट सोपी आहे. पूर्वीप्रमाणे आपण शून्य, उत्तरवर्ती आणि प्रक्षेपण
यांपासून सुरुवात करू आणि संयोजन व आदिम पुनरावर्तन यांखाली बंद करू. फरक
इतकाच की, व्याख्येतील काही पदे परिभाषित नसण्याची शक्यता सामावण्यासाठी
संयोजनाची आणि आदिम पुनरावर्तनाची व्याख्या बदलावी लागते. $f$ आणि $g$ ही
अंशतः फलने असतील, तर $x$ येथे $f$ परिभाषित आहे, म्हणजे $x$ हा $f$ च्या
परिभाषाक्षेत्रात आहे, हे दर्शवण्यासाठी $f(x) \fdefined$ असे लिहू; आणि
याच्या विरुद्ध अर्थासाठी, म्हणजे $x$ येथे $f$ परिभाषित नाही हे दर्शवण्यासाठी
$f(x) \fundefined$ असे लिहू. $f(x)$ आणि $g(x)$ ही दोन्ही अपरिभाषित आहेत,
किंवा दोन्ही परिभाषित असून समान आहेत, हे दर्शवण्यासाठी $f(x) \simeq g(x)$
असे लिहू. अधिक गुंतागुंतीच्या पदांसाठीही ही चिन्हांकने वापरू. $h$ आणि
$g_0$, \dots,~$g_k$ ही सर्व अंशतः फलने असतील, तर
\[h(g_0(\vec x),\dots,g_k(\vec x))
\]
हे पद तेव्हा आणि केवळ तेव्हाच परिभाषित असते, जेव्हा प्रत्येक $g_i$ हे
$\vec x$ येथे परिभाषित असते आणि $h$ हे $g_0(\vec x)$, \dots,~$g_k(\vec x)$
येथे परिभाषित असते, असा संकेत आपण स्वीकारू. हे समजल्यावर अंशतः फलनांसाठीच्या
संयोजनाच्या आणि आदिम पुनरावर्तनाच्या व्याख्या वरीलप्रमाणेच राहतात; फक्त
``$=$'' च्या जागी ``$\simeq$'' लिहावे लागते.

अंशतः फलने मिळवण्यासाठी आदिम पुनरावर्ती फलनांच्या व्याख्येत आपण
\emph{अपरिबद्ध शोध परिकर्मी} जोडू. $f(x,\vec z)$ हे नैसर्गिक संख्यांवरील
कोणतेही अंशतः फलन असेल, तर $\mu x \; f(x,\vec z)$ ची व्याख्या पुढीलप्रमाणे
करू:
\begin{quote}
  $f(0,\vec z), f(1,\vec z), \dots, f(x,\vec
  z)$ ही सर्व मूल्ये परिभाषित असतील आणि $f(x,\vec z) = 0$ असेल असा लघुतम
  $x$, असा $x$ अस्तित्वात असल्यास.
\end{quote}
अन्यथा $\mu x \; f(x,\vec z)$ अपरिभाषित आहे, असे समजायचे. यामुळे
$\mu x \; f(x,\vec z)$ ची एकमेव व्याख्या होते.

\begin{explain}
आपल्या व्याख्येत अल्गोरिदमचा
किंवा कोणत्याही विशिष्ट संगणन-प्रतिमानाचा संदर्भ नाही, हे लक्षात घ्या. पण
संयोजन आणि आदिम पुनरावर्तनाप्रमाणे अपरिबद्ध शोधामागेही एक क्रियात्मक,
संगणकीय सहजकल्पना आहे. अंशतः फलनाच्या संगणनक्षमतेच्या संदर्भात, ज्या
आदानांसाठी फलन अपरिभाषित असते ती आदाने संगणन न थांबणाऱ्या आदानांशी संगत
असतात. $\mu x \; f(x,\vec z)$ संगणित करण्याची प्रक्रिया अशी: शून्य मूल्य
परत येईपर्यंत $f(0,\vec z), f(1,\vec z), f(2,\vec z)$ संगणित करत राहा.
मात्र मधल्या संगणनांपैकी एखादे थांबले नाही, तर $\mu x \; f(x,\vec z)$ चे
संगणनही थांबत नाही.
\end{explain}

$R(x,\vec z)$ हा कोणताही संबंध असेल, तर $\mu x \; R(x,\vec z)$ ची व्याख्या
$\mu x \; (1 \tsub \Char{R}(x,\vec z))$ अशी करतात. दुसऱ्या शब्दांत,
$R(x,\vec z)$ सत्य असलेले $x$ चे लघुतम मूल्य $\mu x \; R(x,\vec z)$ परत
करते. म्हणून $f(x,\vec z)$ हे सर्वत्र परिभाषित फलन असेल, तर
$\mu x \; f(x,\vec z)$ हे $\mu x \; (f(x,\vec z) = 0)$ याच्याशी समान असते.
पण आपली मूळ व्याख्या अधिक सर्वसाधारण आहे, कारण ती $f(x,\vec z)$ सर्वत्र
परिभाषित नसण्याची शक्यता मान्य करते. (याउलट, संबंधाचे जातिबोधक फल नेहमीच
सर्वत्र परिभाषित असते.)

\begin{defn}
नैसर्गिक संख्यांकडून नैसर्गिक संख्यांकडे जाणाऱ्या (विविध स्थानसंख्या असलेल्या)
अंशतः फलनांचा असा लघुतम संच, ज्यात शून्य, उत्तरवर्ती आणि प्रक्षेपणे आहेत आणि
जो संयोजन, आदिम पुनरावर्तन व अपरिबद्ध शोध यांखाली बंद आहे, त्याला
\emph{आंशिक पुनरावर्ती फलनांचा} संच म्हणतात.
\end{defn}

अर्थात, काही आंशिक पुनरावर्ती फलने सर्वत्र परिभाषित असतील, म्हणजे प्रत्येक
आदानासाठी परिभाषित असतील.

\begin{defn}
\label{cmp:rec:par:defn:recursive-fn}
सर्वत्र परिभाषित असलेल्या आंशिक पुनरावर्ती फलनांच्या संचाला
\emph{पुनरावर्ती फलनांचा} संच म्हणतात.
\end{defn}

पुनरावर्ती फलन सर्वत्र परिभाषित आहे यावर भर देण्यासाठी त्याला कधीकधी
``सर्वत्र परिभाषित पुनरावर्ती'' असे म्हणतात.











\section{प्रमाण रूप प्रमेय}\label{cmp:rec:nft:sec}

\begin{thm}[क्लिनीचे प्रमाण रूप प्रमेय]
\label{cmp:rec:nft:thm:kleene-nf}
असा आदिम पुनरावर्ती संबंध $T(e, x, s)$ आणि असे आदिम पुनरावर्ती फलन
$U(s)$ असते की पुढील गुणधर्म पूर्ण होतो: $f$ हे कोणतेही आंशिक पुनरावर्ती
फलन असेल, तर एखाद्या~$e$ साठी सर्व $x$ बाबत
\[f(x) \simeq U(\umin{s}{T(e, x, s)})
\]
असते.
\end{thm}

\begin{explain}
प्रमाण रूप प्रमेयाची सिद्धता गुंतागुंतीची असली, तरी तिची मूलभूत कल्पना
सोपी आहे. प्रत्येक आंशिक पुनरावर्ती फलनाला एक \emph{निर्देशांक}~$e$
असतो; सहज अर्थाने, तो त्या फलनाचा कार्यक्रम किंवा व्याख्या संकेतबद्ध
करणारी संख्या असते. $f(x) \fdefined$ असेल, तर त्या संगणनाची पद्धतशीर
नोंद करून तिला एखाद्या संख्येने~$s$ संकेतबद्ध करता येते. केवळ $x$ आणि
व्याख्या~$e$ वापरून, $s$ ही संख्या आदान~$x$ वरील~$f$ चे संगणन
संकेतबद्ध करते का, हे आदिम पुनरावर्ती पद्धतीने तपासता येते. म्हणून
``निर्देशांक~$e$ असलेल्या फलनाचे आदान~$x$ साठी संगणन आहे आणि $s$
त्या संगणनाचा संकेतांक आहे'' हा संबंध~$T$ आदिम पुनरावर्ती असतो.
संगणनाची पूर्ण नोंद~$s$ दिली असता, $s$ चा परिणाम म्हणजे~$f(x)$ चे मूल्य;
तेही~$s$ पासून आदिम पुनरावर्ती पद्धतीने मिळवता येते.

प्रमाण रूप प्रमेय दाखवते की कोणत्याही आंशिक पुनरावर्ती फलनाच्या
व्याख्येसाठी एकच अपरिबद्ध शोध पुरेसा असतो. म्हणजे, निर्देशांक~$e$
असलेल्या फलनाचे आदान~$x$ वरील संगणन संकेतबद्ध करणारी संख्या सापडेपर्यंत
आपण सर्व संख्यांमध्ये शोध करू शकतो. आंशिक पुनरावर्ती फलनांची ``नावे''
म्हणून या संख्या~$e$ वापरता येतात आणि प्रमेयातील समीकरणाने परिभाषित
केलेल्या~$f$ फलनासाठी $\cfind{e}$ असे लिहिता येते. एका आंशिक
पुनरावर्ती फलनाला एकापेक्षा अधिक निर्देशांक असू शकतात, हे लक्षात घ्या;
खरे तर प्रत्येक आंशिक पुनरावर्ती फलनाला अनंत निर्देशांक असतात.
\end{explain}











\section{थांबण्याची समस्या}\label{cmp:rec:hlt:sec}

सर्वसाधारणपणे \emph{थांबण्याची समस्या} अशी आहे: संगणनक्षम फलनाचे
विनिर्देशन~$e$ (उदा., कार्यक्रम) आणि संख्या~$n$ दिली असता, आदान~$n$
वरील त्या फलनाचे संगणन थांबते का, म्हणजे काही परिणाम देते का, हे
ठरवणे. ही समस्या स्वतः संगणनक्षम फलनाने सोडवता येत नाही, हे ॲलन
ट्यूरिंग यांनी सिद्ध केले. म्हणजे, पुढील फलन संगणनक्षम नाही:
\[h(e, n) =
\begin{cases}
1 & \text{जर संगणन $e$ आदान $n$ वर थांबत असेल तर}\\
0 & \text{अन्यथा,}
\end{cases}
\]

आंशिक पुनरावर्ती फलनांच्या संदर्भात कार्यक्रमाच्या विनिर्देशनाची भूमिका
क्लिनीच्या प्रमाण रूप प्रमेयातील निर्देशांक~$e$ बजावू शकतो. $f$ हे
आंशिक पुनरावर्ती फलन असेल, तर प्रमाण रूप प्रमेयातील समीकरण ज्या
कोणत्याही $e$~साठी खरे असेल, तो~$f$ चा निर्देशांक असतो. संख्या~$e$
दिली असता, प्रमाण रूप प्रमेय सांगते की
\[\cfind{e}(x) \simeq U(\mu s \; T(e, x, s))
\]
हे आंशिक पुनरावर्ती असते; आणि प्रत्येक आंशिक पुनरावर्ती $f\colon \Nat
\to \Nat$ साठी असा $e \in \Nat$ असतो की सर्व~$x \in \Nat$ साठी
$\cfind{e}(x) \simeq f(x)$. खरे तर अशा प्रत्येक $f$ साठी एकच नव्हे,
तर अनंत निर्देशांक~$e$ असतात. \emph{थांबण्याचे फलन}~$h$ पुढीलप्रमाणे
परिभाषित करतात:
\[h(e, x) =
\begin{cases}
1 & \text{जर $\cfind{e}(x) \fdefined$ असेल तर}\\
0 & \text{अन्यथा.}
\end{cases}
\]
$\cfind{e}(x) \fundefined$ असेल तेव्हा $h(e, x) = 0$, हे लक्षात घ्या.
प्रमाण रूपातील सूत्राने प्रत्येक~$e$ साठी आंशिक पुनरावर्ती फलन मिळते;
म्हणून येथे निर्देशांक नसल्याची वेगळी अपवादस्थिती उद्भवत नाही.


\begin{thm}
\label{cmp:rec:hlt:thm:halting-problem}
थांबण्याचे फलन~$h$ आंशिक पुनरावर्ती नाही.
\end{thm}

\begin{proof}
$h$ आंशिक पुनरावर्ती आहे असे समजल्यास, आपण पुढील फलन परिभाषित
करू शकू:
\[d(y) =
\begin{cases}
1 & \text{जर $h(y, y) = 0$ असेल तर}\\
\umin{x}{x \neq x} & \text{अन्यथा.}
\end{cases}
\]
कोणतीही संख्या $x$ ही $x \neq x$ पूर्ण करत नसल्यामुळे $\umin{x}{x \neq
x}$ अस्तित्वात नाही. म्हणून $d(y) \fundefined$ तेव्हा आणि केवळ तेव्हाच,
जेव्हा $h(y,y) \neq 0$. या व्याख्येवरून पुढील गोष्टी मिळतात:
\begin{enumerate}
\item $d(y) \fdefined$ तेव्हा आणि केवळ तेव्हाच, जेव्हा $\cfind{y}(y)
  \fundefined$; प्रत्येक $y$ हा काही आंशिक पुनरावर्ती फलनाचा निर्देशांक
  असल्यामुळे वेगळी अपवादस्थिती नाही.
\item $d(y) \fundefined$ तेव्हा आणि केवळ तेव्हाच, जेव्हा $\cfind{y}(y)
  \fdefined$.
\end{enumerate}
$h$ आंशिक पुनरावर्ती असेल, तर $d$ ही आंशिक पुनरावर्ती असेल. त्यामुळे
क्लिनीच्या प्रमाण रूप प्रमेयानुसार तिला निर्देशांक~$e_d$ असेल.
$h(e_d, e_d)$ चे मूल्य पाहू. $0$ आणि~$1$ असे दोन पर्याय आहेत.
\begin{enumerate}
\item $h(e_d, e_d) = 1$ असेल, तर $\cfind{e_d}(e_d) \fdefined$.
  पण $\cfind{e_d} \simeq d$ आणि $d(e_d)$ हे $h(e_d, e_d) = 0$
  तेव्हा आणि केवळ तेव्हाच परिभाषित असते. म्हणून $h(e_d, e_d) \neq 1$.
\item $h(e_d, e_d) = 0$ असेल, तर $e_d$ हा त्या फलनाचा निर्देशांक असल्याने
  $\cfind{e_d}(e_d) \fundefined$. पण पुन्हा $\cfind{e_d} \simeq d$
  आणि $d(e_d)$ तेव्हा आणि केवळ तेव्हाच अपरिभाषित असते, जेव्हा
  $\cfind{e_d}(e_d) \fdefined$. त्यामुळे या प्रकरणातही विरोधाभास मिळतो.
\end{enumerate}
यावरून $e_d$ हा आंशिक पुनरावर्ती फलनाचा निर्देशांक असताना वरील
दोन्ही पर्याय अशक्य ठरतात. पण $h$ आंशिक पुनरावर्ती असेल, तर $d$ ही
आंशिक पुनरावर्ती असेल आणि तिचा निर्देशांक म्हणून $e_d$ निवडणे योग्य
असेल. म्हणून $h$ आंशिक पुनरावर्ती असू शकत नाही.
\end{proof}












\section{सामान्य पुनरावर्ती फलने}\label{cmp:rec:gen:sec}

सर्वत्र परिभाषित फलनांचा एक संच मिळवण्याचा आणखी एक मार्ग आहे.
सर्वत्र परिभाषित $f(x,\vec z)$ हे फलन \emph{नियमित} असे म्हणू, जर
नैसर्गिक संख्यांच्या प्रत्येक क्रमिका~$\vec z$ साठी असा $x$ असेल की
$f(x,\vec z) = 0$. दुसऱ्या शब्दांत, ज्या फलनांना अपरिबद्ध शोध लागू
केल्यावर सर्वत्र परिभाषित फलन मिळते, तीच नियमित फलने आहेत.
सावधगिरीने अपरिबद्ध शोध नियमित फलनांपुरताच मर्यादित करता येतो:

\begin{defn}
\label{cmp:rec:gen:defn:general-recursive}
\emph{सामान्य पुनरावर्ती फलनांचा} संच हा विविध चलसंख्यांची नैसर्गिक
संख्यांपासून नैसर्गिक संख्यांकडे जाणारी फलने असलेला सर्वांत लहान संच
आहे: त्यात शून्य, उत्तरवर्ती आणि प्रक्षेपणे आहेत, आणि तो संयोजन, आदिम
पुनरावर्तन आणि \emph{नियमित} फलनांना लागू केलेला अपरिबद्ध शोध यांखाली
बंद आहे.
\end{defn}

प्रत्येक सामान्य पुनरावर्ती फलन सर्वत्र परिभाषित आहे, हे स्पष्ट आहे.
\ref{cmp:rec:gen:defn:general-recursive} आणि \ref{cmp:rec:par:defn:recursive-fn}
यांतील फरक असा की दुसऱ्या व्याख्येत मधल्या टप्प्यांवर आंशिक पुनरावर्ती
फलने वापरण्यास परवानगी आहे; शेवटी मिळणारे फलन सर्वत्र परिभाषित असावे,
एवढीच अट आहे. म्हणून ``सामान्य'' हा ऐतिहासिक अवशेष असलेला शब्द
दिशाभूल करणारा आहे: वरवर पाहता \ref{cmp:rec:gen:defn:general-recursive} ही
व्याख्या \ref{cmp:rec:par:defn:recursive-fn} पेक्षा \emph{कमी} सामान्य
आहे. पण सुदैवाने हा फरक केवळ भासमान आहे. व्याख्या वेगळ्या असल्या,
तरी सामान्य पुनरावर्ती फलनांचा संच आणि पुनरावर्ती फलनांचा संच एकच आहे.




